% Generated mechanically from frozen input p12/ko/stacks-geometry.tex.
% Do not edit this generated file; edit the bound locale source instead.

% OK, start here.
%

\setcounter{chapter}{106}
\chapter{대수적 스택의 기하학}



\phantomsection
\label{stacks-geometry-section-phantom}


\section{서론}
\label{stacks-geometry-section-introduction}

\noindent
이 장에서는 대수적 스택의 몇 가지 기하학적 성질을 논한다.
절 \ref{stacks-geometry-section-multiplicities}와
\ref{stacks-geometry-section-dimension-of-algebraic-stacks}의 초판은
Matthew Emerton과 Toby Gee가 작성했으며, 그 원래 형태는
\cite{Emerton-Gee-dim}에서 볼 수 있다.







\section{버설 환}
\label{stacks-geometry-section-versal}

\noindent
이 절에서는 국소 뇌터 기저 위의 유한형 대수적 스택에서
변형환과 국소환 사이의 관계를 밝힌다.

\begin{situation}
\label{stacks-geometry-situation-versal}
여기서 $\mathcal{X}$는 국소 뇌터 스킴 $S$ 위에서
국소적으로 유한형인 대수적 스택이다.
\end{situation}

\noindent
정의는 다음과 같다.

\begin{definition}
\label{stacks-geometry-definition-versal-ring-at-x}
상황 \ref{stacks-geometry-situation-versal}에서 $x_0 : \Spec(k) \to \mathcal{X}$를
사상이라 하자. 여기서 $k$는 $S$ 위의 유한형 체이다.
$\mathcal{X}$의 $x_0$에서의 {\it 버설 환}이란 잉여체가 $k$인
완비 뇌터 국소 $S$-대수 $A$로서, 『Artin의 공리』의 정의
\ref{artin-definition-versal-formal-object}에서와 같은
버설 형식적 대상 $(A, \xi_n, f_n)$이 존재하고
$\xi_1 \cong x_0$인 것($2$-동형사상)을 말한다.
\end{definition}

\noindent
버설 환이 존재하며 매끄러운 인자를 제외하고 유일함을 증명하고자 한다.
이를 위해 『Artin의 공리』의 절
\ref{artin-section-predeformation-categories}에 나오는
전변형 범주를 사용할 것이다. 우리의 상황에서 이들은 항상 변형 범주이다.

\begin{lemma}
\label{stacks-geometry-lemma-deformation-category}
상황 \ref{stacks-geometry-situation-versal}에서 $x_0 : \Spec(k) \to \mathcal{X}$를
사상이라 하자. 여기서 $k$는 $S$ 위의 유한형 체이다.
그러면 $\mathcal{F}_{\mathcal{X}, k, x_0}$는 변형 범주이고,
$T\mathcal{F}_{\mathcal{X}, k, x_0}$와
$\text{Inf}(\mathcal{F}_{\mathcal{X}, k, x_0})$는
유한 차원 $k$-벡터공간이다.
\end{lemma}

\begin{proof}
$\Spec(k) \to S$가 통과하는 아핀 열린집합
$\Spec(\Lambda) \subset S$를 택한다.
『Artin의 공리』의 절 \ref{artin-section-predeformation-categories}에 의해
범주 $\mathcal{C}_\Lambda$ 위의 전변형 범주
$\mathcal{F}_{\mathcal{X}, k, x_0}$를 얻는다.
(앞서 인용한 곳에서 지적했듯이 이 범주는
$\Lambda$의 선택이 아니라 사상 $\Spec(k) \to S$에만 의존한다.)
『Artin의 공리』의 보조정리
\ref{artin-lemma-deformation-category}와
\ref{artin-lemma-algebraic-stack-RS}에 의해
$\mathcal{F}_{\mathcal{X}, k, x_0}$는 실제로 변형 범주이다.
『Artin의 공리』의 보조정리 \ref{artin-lemma-finite-dimension}에 의해
$T\mathcal{F}_{\mathcal{X}, k, x_0}$와
$\text{Inf}(\mathcal{F}_{\mathcal{X}, k, x_0})$가
유한 차원 $k$-벡터공간임을 알 수 있다.
\end{proof}

\begin{lemma}
\label{stacks-geometry-lemma-versal-ring}
상황 \ref{stacks-geometry-situation-versal}에서 $x_0 : \Spec(k) \to \mathcal{X}$를
사상이라 하자. 여기서 $k$는 $S$ 위의 유한형 체이다.
그러면 $\mathcal{X}$의 $x_0$에서의 버설 환이 존재한다.
그러한 두 환 $A$, $A'$가 주어지면, 어떤 $r$에 대해 $S$-대수로서
$A \cong A'[[t_1, \ldots, t_r]]$이거나
$A' \cong A[[t_1, \ldots, t_r]]$이다.
\end{lemma}

\begin{proof}
존재성은 보조정리 \ref{stacks-geometry-lemma-deformation-category}와
『형식적 변형 이론』의 정의
\ref{formal-defos-definition-deformation-category} 및 보조정리
\ref{formal-defos-lemma-RS-implies-S1-S2},
\ref{formal-defos-lemma-versal-object-existence}에서 따른다.
『형식적 변형 이론』의 보조정리
\ref{formal-defos-lemma-minimal-versal}의 유일성 결과에 의해,
$\mathcal{X}$의 $x_0$에서의 ``극소'' 버설 환 $A$가 존재하여
$\mathcal{X}$의 $x_0$에서의 다른 모든 버설 환은 어떤 $r$에 대해
$A[[t_1, \ldots, t_r]]$과 동형이다.
이는 두 번째 명제를 명백히 함의한다.
\end{proof}

\begin{lemma}
\label{stacks-geometry-lemma-versal-ring-field-extension}
상황 \ref{stacks-geometry-situation-versal}에서 $x_0 : \Spec(k) \to \mathcal{X}$를
사상이라 하자. 여기서 $k$는 $S$ 위의 유한형 체이다.
$l/k$를 체의 유한 확대라 하고, 유도된 사상을
$x_{l, 0} : \Spec(l) \to \mathcal{X}$로 나타내자.
$\mathcal{X}$의 $x_0$에서의 버설 환 $A$가 주어지면,
$\mathcal{X}$의 $x_{l, 0}$에서의 버설 환 $A'$가 존재하여
주어진 체의 확대 $l/k$를 유도하고 $\mathfrak m_{A'}$-진 위상에서
형식적으로 매끄러운 $S$-대수 사상 $A \to A'$가 존재한다.
\end{lemma}

\begin{proof}
『Artin의 공리』의 보조정리 \ref{artin-lemma-change-of-field}와
『형식적 변형 이론』의 보조정리
\ref{formal-defos-lemma-change-of-field-versal-ring}에서 곧바로 따른다.
(또한 『Artin의 공리』의 보조정리
\ref{artin-lemma-algebraic-stack-RS}에 의해
$\mathcal{X}$가 (RS)를 만족한다는 사실도 사용한다.)
\end{proof}

\begin{lemma}
\label{stacks-geometry-lemma-compare-versal-ring-completion}
상황 \ref{stacks-geometry-situation-versal}에서 $x : U \to \mathcal{X}$를 사상이라 하자.
여기서 $U$는 $S$ 위에서 국소적으로 유한형인 스킴이다.
$u_0 \in U$를 유한형 점이라 하자.
$k = \kappa(u_0)$로 두고, 유도된 사상을
$x_0 : \Spec(k) \to \mathcal{X}$로 나타내자. 다음은 동치이다.
\begin{enumerate}
\item $x$는 $u_0$에서 버설이다
(『Artin의 공리』의 정의 \ref{artin-definition-versal}),
\item $\hat x : \mathcal{F}_{U, k, u_0} \to \mathcal{F}_{\mathcal{X}, k, x_0}$는
매끄럽다,
\item $x|_{\Spec(\mathcal{O}_{U, u_0}^\wedge)}$에 대응하는 형식적 대상은
버설이다, 그리고
\item $x$의 열린 근방 $U' \subset U$가 존재하여
$x|_{U'} : U' \to \mathcal{X}$가 매끄럽다.
\end{enumerate}
더욱이 이 경우 완비화 $\mathcal{O}_{U, u_0}^\wedge$는
$\mathcal{X}$의 $x_0$에서의 버설 환이다.
\end{lemma}

\begin{proof}
$U \to S$는 (그러한 사상들의 합성으로서) 국소적으로 유한형이므로,
$\Spec(k) \to S$도 (다시 합성으로서) 유한형이다.
따라서 이 명제는 의미가 있다. (1)과 (2)의 동치는
$x$가 $u_0$에서 버설이라는 정의이다.
(1)과 (3)의 동치는 『Artin의 공리』의 보조정리
\ref{artin-lemma-versality-matches}이다. 따라서 (1), (2), (3)은 동치이다.

\medskip\noindent
$x|_{U'}$가 매끄러우면 『Artin의 공리』의 보조정리
\ref{artin-lemma-formally-smooth-on-deformation-categories}에 의해 함자
$\hat x : \mathcal{F}_{U, k, u_0} \to \mathcal{F}_{\mathcal{X}, k, x_0}$는
매끄럽다. 따라서 (4)는 (1), (2), (3)을 함의한다.
역으로 $x$가 $u_0$에서 버설이라고 가정하자.
$V$가 스킴인 전사 매끄러운 사상 $y : V \to \mathcal{X}$를 택한다.
$Z = V \times_\mathcal{X} U$로 두고 $u_0$ 위에 놓이는
유한형 점 $z_0 \in |Z|$를 택한다(이는 『대수공간의 사상』의 보조정리
\ref{spaces-morphisms-lemma-finite-type-points-surjective-morphism}에 의해 가능하다).
『Artin의 공리』의 보조정리 \ref{artin-lemma-base-change-versal}에 의해
사상 $Z \to V$는 $z_0$에서 매끄럽다.
정의에 의해 $z_0$의 열린 근방 $W \subset Z$로서
$W \to V$가 매끄러운 것을 찾을 수 있다. $Z \to U$는 열린 사상이므로,
$U' \subset U$를 $W$의 상으로 두자. 그러면 스택의 매끄러운 사상에 대한
우리의 정의에 의해 $U' \to \mathcal{X}$가 매끄럽다.

\medskip\noindent
마지막 명제는 $\mathcal{O}_{U, u_0}^\wedge$가
$\mathcal{F}_{U, k, u_0}$를 프로표현한다는 사실과 정의에서 따른다.
\end{proof}

\begin{lemma}
\label{stacks-geometry-lemma-characterize-smoothness}
상황 \ref{stacks-geometry-situation-versal}에서 $x_0 : \Spec(k) \to \mathcal{X}$를
사상이라 하자. $\Spec(k) \to S$는 상이 $s$인 유한형 사상이라 하자.
$A$를 $\mathcal{X}$의 $x_0$에서의 버설 환이라 하자. 다음은 동치이다.
\begin{enumerate}
\item $x_0$는 $\mathcal{X} \to S$의 매끄러운 자취에 속한다
(『스택의 사상』의 보조정리 \ref{stacks-morphisms-lemma-where-smooth}),
\item $\mathcal{O}_{S, s} \to A$는 $\mathfrak m_A$-진 위상에서
형식적으로 매끄럽다, 그리고
\item $\mathcal{F}_{\mathcal{X}, k, x_0}$에는 장애가 없다.
\end{enumerate}
\end{lemma}

\begin{proof}
(2)와 (3)의 동치는 『형식적 변형 이론』의 보조정리
\ref{formal-defos-lemma-smooth-power-series-classical}에서 곧바로 따른다.

\medskip\noindent
$\mathcal{O}_{S, s} \to A$가 $\mathfrak m_A$-진 위상에서 형식적으로
매끄러운 것은 $\mathcal{O}_{S, s} \to A' = A[[t_1, \ldots, t_r]]$가
$\mathfrak m_{A'}$-진 위상에서 형식적으로 매끄러운 것과 동치임에 유의하자.
따라서 보조정리 \ref{stacks-geometry-lemma-versal-ring}에 의해 (2)는
버설 환의 선택에 의존하지 않는다.
다음으로 $l/k$를 유한 확대라 하고 보조정리
\ref{stacks-geometry-lemma-versal-ring-field-extension}에서와 같은 $A \to A'$를 택하자.
$\mathcal{O}_{S, s} \to A$가 $\mathfrak m_A$-진 위상에서 형식적으로
매끄러우면, 『대수학 더 보기』의 보조정리
\ref{more-algebra-lemma-compose-formally-smooth}에 의해
$\mathcal{O}_{S, s} \to A'$는 $\mathfrak m_{A'}$-진 위상에서
형식적으로 매끄럽다. 역으로 $\mathcal{O}_{S, s} \to A'$가
$\mathfrak m_{A'}$-진 위상에서 형식적으로 매끄러우면,
$\mathcal{O}_{S, s}^\wedge \to A'$와 $A \to A'$는 정칙이다
(『대수학 더 보기』의 명제 \ref{more-algebra-proposition-fs-regular}).
따라서 $\mathcal{O}_{S, s}^\wedge \to A$도 정칙이고
(『대수학 더 보기』의 보조정리
\ref{more-algebra-lemma-regular-permanence}), 그러므로
$\mathcal{O}_{S, s} \to A$는 $\mathfrak m_A$-진 위상에서
형식적으로 매끄럽다(앞과 같은 보조정리).
따라서 $k$와 $x_0$에 대해 (2)와 (1)이 동치인 것은
$l$과 $x_{0, l}$에 대해 동치인 것과 필요충분하다.

\medskip\noindent
스킴 $U$와 매끄러운 사상 $U \to \mathcal{X}$를
$\Spec(k) \times_\mathcal{X} U$가 공집합이 아니도록 택한다.
유한 확대 $l/k$와 점
$w_0 : \Spec(l) \to \Spec(k) \times_\mathcal{X} U$를 택하자.
$u_0 \in U$를 $w_0$의 상이라 하자.
$l/k$와 $l/\kappa(u_0)$에 위 논의를 적용하면
$u_0$의 경우로 환원할 수 있음을 알 수 있다. 따라서 보조정리
\ref{stacks-geometry-lemma-compare-versal-ring-completion}에 의해
$A = \mathcal{O}_{U, u_0}^\wedge$라 가정해도 좋다.
$x_0$가 $\mathcal{X} \to S$의 매끄러운 자취에 속하는 것은
$u_0$가 $U \to S$의 매끄러운 자취에 속하는 것과 동치임에 유의하자.
예를 들어 『스택의 사상』의 보조정리
\ref{stacks-morphisms-lemma-where-smooth}를 보라.
따라서 (1)과 (2)의 동치는 『대수학 더 보기』의 보조정리
\ref{more-algebra-lemma-formally-smooth-finite-type}에서 따른다.
\end{proof}

\noindent
Artin 근사의 한 귀결을 상기한다.

\begin{lemma}
\label{stacks-geometry-lemma-Artin-approximation-by-smooth-morphism}
상황 \ref{stacks-geometry-situation-versal}에서 $x_0 : \Spec(k) \to \mathcal{X}$를
사상이라 하자. $\Spec(k) \to S$는 상이 $s$인 유한형 사상이라 하자.
$A$를 $\mathcal{X}$의 $x_0$에서의 버설 환이라 하자.
$\mathcal{O}_{S, s}$가 G-환이면, 원천이 스킴인 매끄러운 사상
$U \to \mathcal{X}$와 잉여체가 $k$인 점 $u_0 \in U$를
다음과 같이 찾을 수 있다.
\begin{enumerate}
\item $\Spec(k) \to U \to \mathcal{X}$는 주어진 사상 $x_0$와 일치한다.
\item 동형사상 $\mathcal{O}_{U, u_0}^\wedge \cong A$가 존재한다.
\end{enumerate}
\end{lemma}

\begin{proof}
$(\xi_n, f_n)$을 $A$ 위의 버설 형식적 대상이라 하자.
『Artin의 공리』의 보조정리 \ref{artin-lemma-effective}에 의해
$\xi = (A, \xi_n, f_n)$가 효과적임을 안다.
가정에 의해 $\mathcal{X}$는 $S$ 위에서 국소적으로 유한 표시이고
(『스택의 사상』의 보조정리
\ref{stacks-morphisms-lemma-noetherian-finite-type-finite-presentation}를 사용하라),
따라서 『대수적 스택의 극한』의 명제
\ref{stacks-limits-proposition-characterize-locally-finite-presentation}에 의해
극한을 보존한다. 그러므로 『Artin의 공리』의 보조정리
\ref{artin-lemma-approximate-versal}에서와 같은 Artin 근사에 의해,
유한형 $S$-스킴을 원천으로 하고 잉여체가 $k$인 점 $u_0 \in U$를
포함하는 사상 $U \to \mathcal{X}$를 찾을 수 있다. 이 사상
$U \to \mathcal{X}$는 (1), (2)를
만족하며 $u_0$에서 버설이다. 보조정리
\ref{stacks-geometry-lemma-compare-versal-ring-completion}에 의해 $U$를 축소한 뒤
$U \to \mathcal{X}$가 매끄럽다고 가정해도 좋다.
\end{proof}

\begin{remark}[버설 환의 강화]
\label{stacks-geometry-remark-upgrade}
상황 \ref{stacks-geometry-situation-versal}에서 $x_0 : \Spec(k) \to \mathcal{X}$를
사상이라 하자. 여기서 $k$는 $S$ 위의 유한형 체이다.
$A$를 $\mathcal{X}$의 $x_0$에서의 버설 환이라 하자.
『Artin의 공리』의 보조정리 \ref{artin-lemma-effective}에 의해
우리의 버설 형식적 대상은 사실 $S$ 위의 사상
$$
\Spec(A) \longrightarrow \mathcal{X}
$$
에서 온다. 더욱이 위의 결과들은 각각 이 사상과 양립하도록 강화할 수 있다.
그 목록은 다음과 같다.
\begin{enumerate}
\item 보조정리 \ref{stacks-geometry-lemma-versal-ring}에서 동형사상
$A \cong A'[[t_1, \ldots, t_r]]$ 또는 $A' \cong A[[t_1, \ldots, t_r]]$은
이 사상들과 양립하도록 택할 수 있다.
\item 보조정리 \ref{stacks-geometry-lemma-versal-ring-field-extension}에서 준동형사상
$A \to A'$는 이 사상들과 양립하도록 택할 수 있다.
\item 보조정리 \ref{stacks-geometry-lemma-compare-versal-ring-completion}에서 사상
$\Spec(\mathcal{O}_{U, u_0}^\wedge) \to \mathcal{X}$는 표준 사상
$\Spec(\mathcal{O}_{U, u_0}^\wedge) \to U$와 주어진 사상
$U \to \mathcal{X}$의 합성이다.
\item 보조정리 \ref{stacks-geometry-lemma-Artin-approximation-by-smooth-morphism}에서
동형사상 $\mathcal{O}_{U, u_0}^\wedge \cong A$는
$\Spec(A) \to \mathcal{X}$가 앞 항목의 표준 사상에 대응하도록 택할 수 있다.
\end{enumerate}
각 경우에 명제는 우리의 사상들이 버설 형식적 원소들과 양립한다는 사실에서
따른다. 다만 이때 암묵적으로 나타나는 도식들은 $2$-화살표를
(비표준적으로) 택해야만 $2$-가환한다는 점에 유의하자.
그럼에도 이는 (1)의 암묵적인 사상 $A' \to A$ 또는 $A \to A'$가
형식적 호모토피를 제외하고 잘 정의된다는 뜻이다. 『형식적 변형 이론』의
보조정리 \ref{formal-defos-lemma-versal-unique-up-to-homotopy}를 보라.
\end{remark}

\begin{lemma}
\label{stacks-geometry-lemma-versal-ring-flat}
상황 \ref{stacks-geometry-situation-versal}에서 $x_0 : \Spec(k) \to \mathcal{X}$를
사상이라 하자. 여기서 $k$는 $S$ 위의 유한형 체이다.
$A$를 $\mathcal{X}$의 $x_0$에서의 버설 환이라 하자.
그러면 주석 \ref{stacks-geometry-remark-upgrade}의 사상
$\Spec(A) \to \mathcal{X}$는 평탄하다.
\end{lemma}

\begin{proof}
상의 점에서 $S$의 국소환이 G-환이면, 이는 보조정리
\ref{stacks-geometry-lemma-Artin-approximation-by-smooth-morphism}와
뇌터 국소환에서 그 완비화로 가는 사상이 평탄하다는 사실에서 곧바로 따른다.
일반적인 경우에는 다음과 같이 증명한다.

\medskip\noindent
1단계. $A$와 $A'$가 $\mathcal{X}$의 $x_0$에서의 두 버설 환이면,
결과가 $A$에 대해 참인 것은 $A'$에 대해 참인 것과 동치이다.
실제로 필요하면 $A$와 $A'$를 맞바꾼 뒤, 합성
$$
\Spec(A') \to \Spec(A) \to \mathcal{X}
$$
이 사상 $\Spec(A') \to \mathcal{X}$가 되도록 하는 형식적으로 매끄러운
사상 $\varphi : A \to A'$가 존재한다고 가정해도 좋다.
보조정리 \ref{stacks-geometry-lemma-versal-ring}과 주석 \ref{stacks-geometry-remark-upgrade}를 보라.
$A \to A'$는 충실하게 평탄하므로, 『스택의 사상』의 보조정리
\ref{stacks-morphisms-lemma-composition-flat}과
\ref{stacks-morphisms-lemma-flat-permanence}에서 이 동치를 얻는다.

\medskip\noindent
2단계. $l/k$를 체의 유한 확대라 하자.
$x_{l, 0} : \Spec(l) \to \mathcal{X}$를 유도된 사상이라 하자.
$A$를 $\mathcal{X}$의 $x_0$에서의 버설 환이라 하고, $A \to A'$를
보조정리 \ref{stacks-geometry-lemma-versal-ring-field-extension}에서와 같이 택하자.
그러면 다시 합성
$$
\Spec(A') \to \Spec(A) \to \mathcal{X}
$$
은 사상 $\Spec(A') \to \mathcal{X}$이다. 주석 \ref{stacks-geometry-remark-upgrade}를 보라.
앞과 같이 논하고 1단계를 사용하여 버설 환의 선택은 무관함을 알면,
보조정리가 $x_0$에 대해 성립하는 것은 $x_{l, 0}$에 대해
성립하는 것과 동치임을 알 수 있다.

\medskip\noindent
3단계. 스킴 $U$와 전사 매끄러운 사상 $U \to \mathcal{X}$를 택한다.
그러면 $Z = U \times_\mathcal{X} x_0$ 위의 유한형 점 $z_0$를 택할 수 있다
(이는 공집합이 아닌 대수공간이다). $u_0 \in U$를 $U$에서 $z_0$의 상이라 하자.
$z_0$가 어떤 스킴의 닫힌점 $w_0 \in W$의 상이 되도록 하는
전사 에탈 사상 $W \to Z$를 택한다
(『대수공간의 사상』의 절
\ref{spaces-morphisms-section-points-finite-type}을 보라).
$W \to \Spec(k)$와 $W \to U$가 유한형이므로,
$\kappa(w_0)/k$와 $\kappa(w_0)/\kappa(u_0)$는 체의 유한 확대이다
(『스킴의 사상』의 절 \ref{morphisms-section-points-finite-type}을 보라).
2단계를 두 번 적용하면 $x_0$를 $u_0 \to U \to \mathcal{X}$로
대체할 수 있다. 그러면 우리의 사상은 합성
$$
\Spec(\mathcal{O}_{U, u_0}^\wedge) \to U \to \mathcal{X}
$$
임을 알 수 있다. 첫 번째 화살표는 뇌터 국소환의 완비화가 평탄하므로
평탄하고(『대수학』의 보조정리 \ref{algebra-lemma-completion-flat}),
두 번째 화살표는 매끄러운 사상이 평탄하므로 평탄하다.
평탄성은 합성으로 보존되므로 그 합성도 평탄하다.
\end{proof}

\begin{remark}
\label{stacks-geometry-remark-groupoid-defo}
상황 \ref{stacks-geometry-situation-versal}에서 $x_0 : \Spec(k) \to \mathcal{X}$를
사상이라 하자. 여기서 $k$는 $S$ 위의 유한형 체이다.
보조정리 \ref{stacks-geometry-lemma-deformation-category}와 『형식적 변형 이론』의 정리
\ref{formal-defos-theorem-presentation-deformation-groupoid}에 의해,
$\mathcal{F}_{\mathcal{X}, k, x_0}$가 범주 $\mathcal{C}_\Lambda$ 위의
함자들로 이루어진 매끄러운 프로표현가능 준군에 의한 표시를 가짐을 안다.
정의를 풀어 쓰면 다음을 택할 수 있다는 뜻이다.
\begin{enumerate}
\item 잉여체가 $k$인 뇌터 완비 국소 $\Lambda$-대수 $A$와
$A$ 위의 $\mathcal{F}_{\mathcal{X}, k, x_0}$의 버설 형식적 대상 $\xi$,
\item 잉여체가 $k$인 뇌터 완비 국소 $\Lambda$-대수 $B$와 동형사상
$$
\underline{B}|_{\mathcal{C}_\Lambda}
\longrightarrow
\underline{A}|_{\mathcal{C}_\Lambda}
\times_{\underline{\xi}, \mathcal{F}_{\mathcal{X}, k, x_0}, \underline{\xi}}
\underline{A}|_{\mathcal{C}_\Lambda}
$$
\end{enumerate}
사영들은 형식적으로 매끄러운 사상 $t : A \to B$와
$s : A \to B$에 대응한다($\xi$가 버설이기 때문이다).
사상 $c : B \to B \widehat{\otimes}_{s, A, t} B$가 존재하여
$(A, B, s, t, c)$를 잉여체가 $k$인 뇌터 완비 국소
$\Lambda$-대수들의 범주에서 쌍대준군으로 만든다
(프로표현가능 함자들에 대해서는 이 사상을 『형식적 변형 이론』의
보조정리 \ref{formal-defos-lemma-presentation-construction}에서 구성한다).
마지막으로, 인용한 정리에 따르면 $\xi$는 $\mathcal{C}_\Lambda$ 위에서
쌍대섬유화된 준군들의 동치
$$
[\underline{A}|_{\mathcal{C}_\Lambda} / \underline{B}|_{\mathcal{C}_\Lambda}]
\longrightarrow
\mathcal{F}_{\mathcal{X}, k, x_0}
$$
를 유도한다. 사실 완비화된 범주 $\widehat{\mathcal{C}}_\Lambda$ 위에서
쌍대섬유화된 준군들의 동치
$$
[\underline{A}/\underline{B}]
\longrightarrow
\widehat{\mathcal{F}}_{\mathcal{X}, k, x_0}
$$
도 얻는다(이것이 작동하는 이유는 『형식적 변형 이론』의 절
\ref{formal-defos-section-prorepresentable-groupoids-in-functors}에 있는
논의를 보라). 물론 $A$는 $\mathcal{X}$의 $x_0$에서의 버설 환이다.
\end{remark}








\section{대수적 스택의 성분의 중복도}
\label{stacks-geometry-section-multiplicities}

\noindent
$X$가 국소 뇌터 스킴이면, $X$를(단순히 위상공간으로 생각하여)
기약 성분들의 합집합 $X = \bigcup T_i$로 쓸 수 있다.
각 기약 성분은 유일한 일반점 $\xi_i$의 폐포이고,
국소환 $\mathcal O_{X,\xi_i}$는 국소 아르틴 환이다.
{\it $T_i$를 따른 $X$의 중복도} 또는
{\it $X$에서 $T_i$의 중복도}를
$$
m_{T_i, X} = \text{length}_{\mathcal O_{X, \xi_i}} \mathcal O_{X, \xi_i}
$$
로 정의할 수 있다. 다시 말해 이는 국소 아르틴 환의 길이이다.
『Chow 호몰로지』의 절
\ref{chow-section-cycle-of-closed-subscheme}과 비교하라.

\medskip\noindent
여기서 우리의 목표는 이 정의를 국소 뇌터 대수적 스택으로 일반화하는 것이다.
$\mathcal{X}$가 스택이면 그 위상공간 $|\mathcal{X}|$는
(『스택의 성질』의 정의
\ref{stacks-properties-definition-topological-space}를 보라)
국소 뇌터이다(『스택의 사상』의 보조정리
\ref{stacks-morphisms-lemma-Noetherian-topology}).
$|\mathcal{X}|$의 기약 성분을 때로 $\mathcal{X}$의 기약 성분이라고 부른다.
$\mathcal{X}$가 준분리이면 $|\mathcal{X}|$는 소버 공간이지만
(『스택의 사상』의 보조정리 \ref{stacks-morphisms-lemma-sober-qs}),
준분리가 아닌 경우에는 그렇지 않을 수 있다. 예를 들어 준분리가 아닌
대수공간 $X = \mathbf{A}^1_\mathbf{C}/\mathbf{Z}$을 생각하라.
더욱이 $|\mathcal{X}|$ 위에는 그 줄기를 사용하여 중복도를 정의할 수 있는
구조층이 없다.

\begin{lemma}
\label{stacks-geometry-lemma-map-of-components}
$f : U \to \mathcal{X}$를 스킴에서 국소 뇌터 대수적 스택으로 가는
매끄러운 사상이라 하자. $|U|$의 임의의 기약 성분의 상의 폐포는
$|\mathcal{X}|$의 기약 성분이다. $U \to \mathcal{X}$가 전사이면,
$|\mathcal{X}|$의 모든 기약 성분을 이 방식으로 얻는다.
\end{lemma}

\begin{proof}
『스택의 성질』의 보조정리
\ref{stacks-properties-lemma-topology-points}에 의해
사상 $|U| \to |\mathcal{X}|$는 연속이고 열린 사상이다.
$T \subset |U|$를 기약 성분이라 하자. $U$는 국소 뇌터이므로
$T$에 포함되는 공집합이 아닌 아핀 열린집합 $W \subset U$를 찾을 수 있다.
그러면 $f(T) \subset |\mathcal{X}|$는 기약이고 공집합이 아닌 열린부분집합
$f(W)$를 포함한다. 따라서 $f(T)$의 폐포는 기약이고 공집합이 아닌
열린집합을 포함한다. 그러므로 이 폐포는 기약 성분이다.

\medskip\noindent
$U \to \mathcal{X}$가 전사라고 가정하고,
$Z \subset |\mathcal{X}|$를 기약 성분이라 하자.
$Z$와 만나는 $|\mathcal{X}|$의 뇌터 열린부분집합 $V$를 택한다.
$V$에서 다른 기약 성분들을 제거한 뒤 $V \subset Z$라 가정해도 좋다.
공집합이 아닌 열린집합 $f^{-1}(V) \subset |U|$의 기약 성분 하나를 택하고,
$T \subset |U|$를 그 폐포라 하자. 이는 $|U|$의 기약 성분이며,
$T$의 선택에 의해 $f(T)$의 폐포는 $Z$와 같아야 한다.
\end{proof}

\noindent
앞의 보조정리는 특히 국소 뇌터 스킴 사이의 매끄러운 사상에도 적용된다.
다음 보조정리의 명제에서는 이 특수한 경우를 암묵적으로 사용한다.

\begin{lemma}
\label{stacks-geometry-lemma-multiplicities}
$U \to X$를 국소 뇌터 스킴들의 매끄러운 사상이라 하자.
$T'$가 $U$의 기약 성분이라 하자. $T'$의 상의 폐포로 얻어지는
$X$의 기약 성분을 $T$라 하자. 그러면 $m_{T', U} = m_{T, X}$이다.
\end{lemma}

\begin{proof}
$T'$의 일반점을 $\xi'$, $T$의 일반점을 $\xi$라 쓰자.
$A = \mathcal{O}_{X, \xi}$ 및 $B = \mathcal{O}_{U, \xi'}$로 둔다.
$\text{length}_A A = \text{length}_B B$임을 보여야 한다.
$A \to B$는 환의 평탄 국소 준동형사상이므로
(매끄러운 사상은 평탄하기 때문이다), 『대수학』의 보조정리
\ref{algebra-lemma-pullback-module}에 의해
$$
\text{length}_A(A) \text{length}_B(B/\mathfrak m_A B) =
\text{length}_B(B)
$$
이다. 따라서 $\mathfrak m_A B = \mathfrak m_B$, 또는 동치로
$B/\mathfrak m_A B$가 축소됨을 보이면 충분하다.
$U \to X$가 매끄러우므로 그 밑변환
$U_{\xi} \to \Spec \kappa(\xi)$도 매끄럽다.
$U_{\xi}$는 체 위의 매끄러운 스킴이므로 축소되고,
따라서 임의의 점에서 그 국소환도 축소된다
(『다양체』의 보조정리
\ref{varieties-lemma-smooth-geometrically-normal}). 특히
$$
B/\mathfrak m_A B =
\mathcal{O}_{U, \xi'}/\mathfrak m_{X, \xi}\mathcal{O}_{U, \xi'} =
\mathcal{O}_{U_\xi, \xi'}
$$
는 축소된다. 이것이 요구한 바이다.
\end{proof}

\noindent
이 결과를 사용하면 매끄러운 국소 관점으로 살펴봄으로써
중복도의 좋은 개념이 존재함을 보일 수 있다.

\begin{lemma}
\label{stacks-geometry-lemma-multiplicity}
$U_1 \to \mathcal{X}$와 $U_2 \to \mathcal{X}$를 스킴들에서
국소 뇌터 대수적 스택 $\mathcal{X}$로 가는 두 매끄러운 사상이라 하자.
$T_1'$과 $T_2'$를 각각 $|U_1|$과 $|U_2|$의 기약 성분이라 하자.
$T_1'$과 $T_2'$의 상들의 폐포가 $|\mathcal{X}|$의 같은 기약 성분 $T$라고
가정하자. 그러면 $m_{T_1', U_1} = m_{T_2', U_2}$이다.
\end{lemma}

\begin{proof}
$V_1$과 $V_2$를 각각 $T_1'$과 $T'_2$의 조밀한 부분집합으로서,
각각 $U_1$과 $U_2$에서 열린 것으로 택한다
(보조정리 \ref{stacks-geometry-lemma-map-of-components}의 증명을 보라).
$|V_1|$과 $|V_2|$의 $|\mathcal{X}|$ 안에서의 상은 기약 부분집합 $T$의
공집합이 아닌 열린부분집합들이므로, 그 교집합은 공집합이 아니다.
『스택의 성질』의 보조정리
\ref{stacks-properties-lemma-points-cartesian}에 의해 사상
$|V_1 \times_\mathcal{X} V_2| \to |V_1| \times_{|\mathcal{X}|} |V_2|$는
전사이다. 따라서 $V_1 \times_\mathcal{X} V_2$는 공집합이 아닌
대수공간이다. 그러므로 원천이 (공집합이 아닌) 스킴인 에탈 전사
$V \to V_1 \times_\mathcal{X} V_2$를 택할 수 있다.
$T'$를 $V$의 임의의 기약 성분이라 하면, 보조정리
\ref{stacks-geometry-lemma-map-of-components}에 의해 $T'$의 $U_1$(각각 $U_2$) 안에서의
상의 폐포는 $T'_1$(각각 $T'_2$)과 같다.

\medskip\noindent
보조정리 \ref{stacks-geometry-lemma-multiplicities}를 두 번 적용하면 요구한 대로
$$
m_{T_1', U_1} = m_{T', V} = m_{T_2', U_2},
$$
를 얻는다.
\end{proof}

\begin{definition}
\label{stacks-geometry-definition-multiplicity}
$\mathcal{X}$를 국소 뇌터 대수적 스택이라 하자.
$T \subset |\mathcal{X}|$를 기약 성분이라 하자.
$\mathcal{X}$에서 $T$의 {\it 중복도}를
$m_{T, \mathcal{X}} = m_{T', U}$로 정의한다. 여기서
$f : U \to \mathcal{X}$는 스킴에서 오는 매끄러운 사상이고,
$T' \subset |U|$는 $f(T') \subset T$인 기약 성분이다.
\end{definition}

\noindent
보조정리 \ref{stacks-geometry-lemma-map-of-components}와 \ref{stacks-geometry-lemma-multiplicity}에 의해,
이는 $f : U \to \mathcal{X}$의 선택과 $T$로 사상되는
기약 성분 $T'$의 선택에 무관하다.

\medskip\noindent
마지막으로, $\mathcal{X}$의 기약 성분을 닫힌 부분스택으로 생각하는 것이
때때로 편리함을 언급한다. 이를 위해 $T \subset |\mathcal{X}|$가
기약 성분이면 $|\mathcal{T}| = T$인 유일한 축소 닫힌 부분스택
$\mathcal{T} \subset \mathcal{X}$를 생각할 수 있다. 『스택의 성질』의 정의
\ref{stacks-properties-definition-reduced-induced-stack}를 보라.
$\mathcal{X}$가 준분리이면 기약 성분은 정수적 스택이다.
더 자세한 논의는 『스택의 사상』의 절
\ref{stacks-morphisms-section-integral-stacks}을 보라.



\section{형식적 가지와 중복도}
\label{stacks-geometry-section-formal-branches}

\noindent
정의 \ref{stacks-geometry-definition-multiplicity}에서 주어진 기약 성분의 중복도 개념과,
유한형 점에서 $\mathcal{X}$의 버설 환(의 스펙트럼)의 기약 성분들의
중복도라는 관련 개념을 비교해 두면 편리할 것이다.

\medskip\noindent
상황 \ref{stacks-geometry-situation-versal}에서 $x_0 : \Spec(k) \to \mathcal{X}$를
사상이라 하자. 여기서 $k$는 $S$ 위의 유한형 체이다.
$A$, $A'$를 $\mathcal{X}$의 $x_0$에서의 버설 환이라 하자.
필요하면 $A$와 $A'$를 맞바꾼 뒤, 버설 형식적 대상들과 양립하는
형식적으로 매끄러운\footnote{$A'$가 $A$ 위의 멱급수환과
동형이 된다는 의미이다.} 사상 $\varphi : A \to A'$가 존재함을 안다.
보조정리 \ref{stacks-geometry-lemma-versal-ring}과 주석 \ref{stacks-geometry-remark-upgrade}를 보라.
더욱이 $\varphi$는 형식적 호모토피를 제외하고 잘 정의된다.
『형식적 변형 이론』의 보조정리
\ref{formal-defos-lemma-versal-unique-up-to-homotopy}를 보라.
특히 『형식적 변형 이론』의 보조정리
\ref{formal-defos-lemma-homotopic-minimal-prime}에 의해
$\varphi(\mathfrak p)A'$는 $A'$의 잘 정의된 아이디얼이다.
$A \to A'$는 형식적으로 매끄러우므로, 실제로
$\varphi(\mathfrak p)A'$는 $A'$의 극소 소 아이디얼이고,
$A'$의 모든 극소 소 아이디얼은 유일한 극소 소 아이디얼
$\mathfrak p \subset A$에 대해 이 형태이다
(이 모든 것은 $A'$를 $A$ 위의 멱급수환으로 써서 쉽게 증명할 수 있다).
따라서 극소 소 아이디얼이 기약 성분에 대응한다는 사실을 상기하면
다음 정의는 의미가 있다.

\begin{definition}
\label{stacks-geometry-definition-formal-branches}
$\mathcal{X}$를 국소 뇌터 스킴 $S$ 위에서 국소적으로 유한형인
대수적 스택이라 하자. $x_0 : \Spec(k) \to \mathcal{X}$를 사상이라 하자.
여기서 $k$는 $S$ 위의 유한형 체이다.
{\it $x_0$를 지나는 $\mathcal{X}$의 형식적 가지}란
$\mathcal{X}$의 $x_0$에서의 버설 환을 임의로 택했을 때
$\Spec(A)$의 기약 성분들의 집합을 말하며, $A$의 선택이 다를 때에는
위에서 설명한 절차로 이 집합들을 동일시한다.
\end{definition}

\noindent
정의 \ref{stacks-geometry-definition-formal-branches}의 상황에서 유한 확대 $l/k$가
주어졌다고 하자. $x_{l, 0} : \Spec(l) \to \mathcal{X}$를
$\Spec(l) \to \Spec(k)$와 $x_0$의 합성으로 두자.
$A \to A'$를 보조정리 \ref{stacks-geometry-lemma-versal-ring-field-extension}에서와 같이 택하자.
$A \to A'$가 충실하게 평탄하므로 사상
$$
\Spec(A') \to \Spec(A)
$$
은 기약 성분의 (일반점)을 기약 성분의 (일반점)으로 보낸다.
이는 전사 사상이지만 일반적으로 전단사는 아니다.
다시 말해 전사 사상
$$
\text{형식적 가지: }\mathcal{X}\text{의 점 }x_{l, 0}
\longrightarrow
\text{형식적 가지: }\mathcal{X}\text{의 점 }x_0
$$
를 얻는다. $l/k$가 순수 비분리이면 이 사상은 단사이기도 하다
(이를 필요로 하게 되면 여기에 정확한 명제와 증명을 덧붙일 것이다).

\begin{lemma}
\label{stacks-geometry-lemma-branches}
정의 \ref{stacks-geometry-definition-formal-branches}의 상황에서,
$x_0$를 지나는 $\mathcal{X}$의 형식적 가지들의 집합에서
$|\mathcal{X}|$ 안에서 $x_0$를 포함하는 $|\mathcal{X}|$의
기약 성분들의 집합으로 가는 표준 전사가 존재한다.
\end{lemma}

\begin{proof}
$A$를 정의 \ref{stacks-geometry-definition-formal-branches}에서와 같이 택하고,
$\Spec(A) \to \mathcal{X}$를 주석 \ref{stacks-geometry-remark-upgrade}에서와 같이 택하자.
$\Spec(A)$의 기약 성분의 일반점은 $|\mathcal{X}|$의
기약 성분의 일반점으로 사상됨을 주장한다.
스킴 $U$와 전사 매끄러운 사상 $U \to \mathcal{X}$를 택하자.
도식
$$
\xymatrix{
\Spec(A) \times_\mathcal{X} U \ar[d]_p \ar[r]_-q & U \ar[d]^f \\
\Spec(A) \ar[r]^j & \mathcal{X}
}
$$
을 생각하자. 보조정리 \ref{stacks-geometry-lemma-versal-ring-flat}에 의해 $j$는 평탄하고,
따라서 $q$도 평탄하다. 한편 $f$는 전사 매끄러운 사상이므로
$p$도 전사 매끄러운 사상이다. 이는 기약 성분의 임의의 일반점
$\eta \in \Spec(A)$가 대수공간
$\Spec(A) \times_\mathcal{X} U$의 여차원 $0$인 점 $\eta'$의 상임을
함의한다(표기법과 에탈 국소환에 대한 내려가기를 사용하려면
『대수공간의 성질』의 절
\ref{spaces-properties-section-generic-points}을 보라).
$q$가 평탄하므로 $q(\eta')$는 $U$의 여차원 $0$인 점이다(같은 논증).
$U$는 스킴이므로 $q(\eta')$는 $U$의 한 기약 성분의 일반점이다.
따라서 주장한 대로, 보조정리 \ref{stacks-geometry-lemma-map-of-components}에 의해
$|\mathcal{X}|$ 안에서 $q(\eta')$의 상의 폐포는 기약 성분이다.

\medskip\noindent
이 주장은 원하는 사상을 정의하는 방법을 명백히 제공한다.
그것이 전사임을 보이기 위해 $|\mathcal{X}|$에서 $x_0$로 사상되는
$u_0 \in U$를 택하자. $u_0$의 아핀 열린 근방 $U' \subset U$를 택한다.
$U'$를 축소한 뒤 $U'$의 모든 기약 성분이 $u_0$를 지난다고 가정해도 좋다.
그런 다음 $\mathcal{X}$를 $|U'| \to |\mathcal{X}|$의 상에 대응하는
열린 부분스택으로 대체할 수 있다. 따라서 $U$가 아핀이고,
$x_0 \in |\mathcal{X}|$로 사상되는 점 $u_0$를 가지며,
$U$의 모든 기약 성분이 $u_0$를 지난다고 가정해도 좋다.
『스택의 성질』의 보조정리
\ref{stacks-properties-lemma-points-cartesian}에 의해
$\Spec(A)$의 닫힌점과 $u_0$로 사상되는 점
$t \in |\Spec(A) \times_\mathcal{X} U|$가 존재한다.
평탄 국소환 준동형사상들
$$
A \longrightarrow
\mathcal{O}_{\Spec(A) \times_\mathcal{X} U, \overline{t}}
\longleftarrow \mathcal{O}_{U, u_0}
$$
에 내려가기를 사용하면, $\mathcal{O}_{U, u_0}$의 모든 극소 소 아이디얼은
가운데 국소환의 어떤 극소 소 아이디얼의 상이고, 그러한 극소 소 아이디얼은
$A$의 극소 소 아이디얼로 사상됨을 알 수 있다. 이는 전사성을 증명한다.
몇 가지 세부사항은 생략한다.
\end{proof}

\noindent
$A$를 뇌터 완비 국소환이라 하자. 그러면 $\Spec(A)$의 기약 성분들은
중복도를 가진다. 절 \ref{stacks-geometry-section-multiplicities}의 서론을 보라.
$A' = A[[t_1, \ldots, t_r]]$이면 사상
$\Spec(A') \to \Spec(A)$는 기약 성분들 사이에서 중복도를 보존하는
전단사를 유도한다(쉬운 증명은 생략한다).
이 사실과 정의 \ref{stacks-geometry-definition-formal-branches} 앞의 논의에 의해
다음 정의는 의미가 있다.

\begin{definition}
\label{stacks-geometry-definition-multiplicity-formal-branches}
$\mathcal{X}$를 국소 뇌터 스킴 $S$ 위에서 국소적으로 유한형인
대수적 스택이라 하자. $x_0 : \Spec(k) \to \mathcal{X}$를 사상이라 하자.
여기서 $k$는 $S$ 위의 유한형 체이다.
{\it $x_0$를 지나는 $\mathcal{X}$의 형식적 가지의 중복도}란
$\mathcal{X}$의 $x_0$에서의 버설 환을 임의로 택했을 때
$\Spec(A)$의 대응하는 기약 성분의 중복도를 말한다(위 논의를 보라).
\end{definition}

\begin{lemma}
\label{stacks-geometry-lemma-branches-multiplicity}
$\mathcal{X}$를 국소 뇌터 스킴 $S$ 위에서 국소적으로 유한형인
대수적 스택이라 하자. $x_0 : \Spec(k) \to \mathcal{X}$를 사상이라 하자.
여기서 $k$는 $S$ 위의 유한형 체이고 그 상은 $s \in S$이다.
$\mathcal{O}_{S, s}$가 G-환이면, 보조정리 \ref{stacks-geometry-lemma-branches}의 사상은
중복도를 보존한다.
\end{lemma}

\begin{proof}
보조정리 \ref{stacks-geometry-lemma-Artin-approximation-by-smooth-morphism}에 의해,
$U$가 스킴인 매끄러운 사상 $U \to \mathcal{X}$와 $U$의 $k$-값 점
$u_0$가 존재하고 $\mathcal{O}_{U, u_0}^\wedge$가 $\mathcal{X}$의
$x_0$에서의 버설 환이라고 가정해도 좋다.
보조정리 \ref{stacks-geometry-lemma-branches}의 증명에서 우리의 사상을 구성한 방식에 의해
($A = \mathcal{O}_{U, u_0}^\wedge$이므로 크게 단순화된다),
다음을 보이면 충분하다. $u_0$를 지나는 $U$의 기약 성분의 중복도는
그 성분으로 사상되는 $\Spec(\mathcal{O}_{U, u_0}^\wedge)$의
임의의 기약 성분의 중복도와 같다.

\medskip\noindent
가환대수로 옮기면 다음과 같다. $C = \mathcal{O}_{U, u_0}$로 두자.
이는 $\mathcal{O}_{S, s}$ 위에서 본질적으로 유한형이므로 G-환이다
(『대수학 더 보기』의 명제
\ref{more-algebra-proposition-finite-type-over-G-ring}).
또 $A = C^\wedge$이다. 따라서 $C \to A$는 정칙 환 사상이다.
$\mathfrak q \subset C$를 극소 소 아이디얼이라 하고,
$\mathfrak p \subset A$를 $\mathfrak q$ 위에 놓이는 극소 소 아이디얼이라 하자.
그러면
$$
R = C_\mathfrak p \longrightarrow A_\mathfrak p = R'
$$
은 아르틴 국소환들의 정칙 환 사상이다. 그러한 환 사상에 대해서는 항상
$$
\text{length}_R R = \text{length}_{R'} R'
$$
이다. 좌변은 $U$ 위에서 우리 성분의 중복도이고 우변은
$\Spec(A)$ 위에서 우리 성분의 중복도이므로, 이것이 보여야 할 바이다.
등식을 보이기 위해 먼저 『대수학』의 보조정리
\ref{algebra-lemma-pullback-module}에 의해
$$
\text{length}_R(R) \text{length}_{R'}(R'/\mathfrak m_R R') =
\text{length}_{R'}(R')
$$
임을 사용한다. 따라서 $\mathfrak m_R R' = \mathfrak m_{R'}$임을 보이면
충분한데, 이는 0차원 국소환들의 정칙 준동형사상이라는 사실의 귀결이다.
\end{proof}













\section{대수적 스택의 차원론}
\label{stacks-geometry-section-dimension-of-algebraic-stacks}

\noindent
우리가 아는 한, 문헌에서 대수적 스택의 차원론에 관한 주요 결과는
여차원과 상대 차원의 개념을 연구한 \cite{Osserman}의 결과이다.
우리는 한 점에서 대수적 스택의 차원이라는 개념을 더 자세히 살펴보고,
사상의 원천에 있는 한 점에서 섬유의 차원을 그 원천과 표적의 차원에
연결하는 여러 결과를 증명한다. 또한 (적절한 가정 아래에서)
한 점에서 대수적 스택의 차원을 버설 환으로 계산할 수 있게 하는 결과인
아래의 보조정리 \ref{stacks-geometry-lemma-dimension-formula}도 증명한다.

\medskip\noindent
항상 결과를 최적으로 만들려고 하지는 않았지만, 불필요한 가정은 대체로
피하려고 했다. 그러나 한 점에서 대수적 스택의 어떤 성질과 그 점에서
이 스택의 버설 환의 성질을 비교하는 일부 결과에서는, 모든 국소환이
$G$-환인 국소 뇌터 스킴을 기저로 하여 그 위에서 국소적으로 유한 표시인
대수적 스택으로 범위를 제한했다. 이렇게 하면 Artin 근사를 이용하여
버설 환의 기하학과 스택 자체의 기하학을 편리하게 비교할 수 있다.
다만 우리가 증명하는 여러 명제가 모두 참이기 위해 이 제한적인 가정이
필요하지는 않을 수도 있다. 그렇지만 염두에 둔 응용에서는 이 가정이
만족되므로, 도움이 될 때에는 기꺼이 이를 가정했다.

\medskip\noindent
$X$가 스킴이면 $X$의 차원 $\dim(X)$을 $X$의 바탕 위상공간의
크룰 차원으로 정의한다. 한편 $x$가 $X$의 점이면 $x$에서 $X$의 차원
$\dim_x (X)$을 $x$를 포함하는 $X$의 열린부분집합 $U$들의 차원의
최솟값으로 정의한다. 『스킴의 성질』의 정의
\ref{properties-definition-dimension}을 보라.
『스킴의 성질』의 보조정리 \ref{properties-lemma-dimension}에 의해
$\dim(X) = \sup_{x \in X} \dim_x(X)$라는 관계가 성립한다.
$X$가 국소 뇌터이면 $\dim_x(X)$는 $x$를 지나는 $X$의 기약 성분들의
$x$에서의 차원의 상한과 일치한다.

\medskip\noindent
$X$가 대수공간이고 $x \in |X|$이면 $\dim_x X = \dim_u U$로 정의한다.
여기서 $U$는 에탈 전사 $U \to X$를 갖는 임의의 스킴이고,
$u\in U$는 $x$ 위에 놓이는 임의의 점이다. 『대수공간의 성질』의 정의
\ref{spaces-properties-definition-dimension-at-point}를 보라.
또 $\dim(X) = \sup_{x \in |X|} \dim_x(X)$로 둔다.
『대수공간의 성질』의 정의
\ref{spaces-properties-definition-dimension}을 보라.

\begin{remark}
\label{stacks-geometry-remark-dimension-algebraic-space}
일반적으로 한 점 $x$에서 대수공간 $X$의 차원은 $x$에서 바탕 위상공간
$|X|$의 차원과 일치하지 않을 수 있다. 예를 들어 $k$가 표수 0인 체이고
$X =  \mathbf{A}^1_k / \mathbf{Z}$이면, $X$는 각 점에서 차원이
$1$($\mathbf{A}^1_k$의 차원)이지만 $|X|$는 비이산 위상을 가지므로
크룰 차원이 0이다. 한편 『대수공간』의 예
\ref{spaces-example-infinite-product}에는 각 점에서 차원이 $0$이지만
$|X|$가 크룰 차원 $1$인 기약공간이고 일반점을 갖는 대수공간의 예가 있다
(따라서 $|X|$은 모든 점에서 차원이 $1$이다). 이 예에 관한 논의는
『대수공간의 성질』의 절
\ref{spaces-properties-section-dimension}도 보라.

\medskip\noindent
한편 $X$가 『적절한 대수공간』의 정의
\ref{decent-spaces-definition-very-reasonable}의 의미에서
{\it 적절한} 대수공간이면(특히 $X$가 준분리이면 그렇다.
『적절한 대수공간』의 절
\ref{decent-spaces-section-reasonable-decent}을 보라), 실제로
$x$에서 $X$의 차원은 $x$에서 $|X|$의 차원과 일치한다.
『적절한 대수공간』의 보조정리
\ref{decent-spaces-lemma-dimension-decent-space}를 보라.
\end{remark}

\noindent
대수적 스택의 차원을 정의하기 위해서는 먼저 원천이 대수공간이고
표적이 대수적 스택인 사상의 원천의 한 점에서 상대 차원이라는 개념을
도입하는 것이 유용하다. 정의가 조금 복잡한 까닭은 (스킴의 경우와 달리)
대수적 스택이나 대수공간의 점을 체의 스펙트럼에서 오는 사상으로
나타낼 수 없고 오직 그러한 사상들의 동치류로만 나타낼 수 있기 때문이다.

\begin{definition}
\label{stacks-geometry-definition-relative-dimension}
$f : T \to \mathcal{X}$가 대수공간에서 대수적 스택으로 가는
국소적으로 유한형인 사상이고, $t \in |T|$가 상이
$x \in | \mathcal{X}|$인 점이면, $t$에서 $f$의 {\it 상대 차원}을
$\dim_t(T_x)$로 표기하고 다음과 같이 정의한다.
$x$를 나타내는, 체의 스펙트럼을 원천으로 하는 사상
$\Spec k \to \mathcal{X}$를 택하고, $|T|$로의 사영 아래에서 $t$로 사상되는
점 $t' \in |T \times_{\mathcal{X}} \Spec k|$를 택한다
(그러한 점 $t'$는 『스택의 성질』의 보조정리
\ref{stacks-properties-lemma-points-cartesian}에 의해 존재한다). 그러면
$$
\dim_t(T_x) = \dim_{t'}(T \times_{\mathcal{X}} \Spec k ).
$$
\end{definition}

\noindent
$T$는 대수공간이고 $\mathcal{X}$는 대수적 스택이므로 섬유곱
$T \times_{\mathcal{X}} \Spec k$는 대수공간이고, 따라서 이 정의에서
제안한 우변의 양은 실제로 정의되어 있음에 유의하자(위 논의를 보라).

\begin{remark}
\label{stacks-geometry-remark-relative-dimension}
(1)
예를 들어 스킴들의 국소적으로 유한형인 사상의 상대 차원이 밑변환 아래에서
불변임을 사용하면(『스킴의 사상』의 보조정리
\ref{morphisms-lemma-dimension-fibre-after-base-change}를 보라),
$\dim_t(T_x)$가 그 계산에 사용한 선택과 무관하게 잘 정의됨을 쉽게 확인한다.

\medskip\noindent
(2)
$\mathcal{X}$도 대수공간인 경우, 이 정의가 『대수공간의 사상』의 정의
\ref{spaces-morphisms-definition-dimension-fibre}에서 주어진 상대 차원의
정의와 일치함을 곧바로 확인할 수 있다.
\end{remark}

\noindent
이제 국소 뇌터 대수적 스택의 차원 연구의 토대가 되는
다음 보조정리를 상기한다.

\begin{lemma}
\label{stacks-geometry-lemma-behaviour-of-dimensions-wrt-smooth-morphisms}
$f: U \to X$가 국소 뇌터 대수공간들의 매끄러운 사상이고,
$u \in |U|$의 상이 $x \in |X|$이면
$$
\dim_u (U) = \dim_x(X) + \dim_{u} (U_x)
$$
이다. 여기서 $\dim_u (U_x)$는 정의
\ref{stacks-geometry-definition-relative-dimension}에 의해 정의한다.
\end{lemma}

\begin{proof}
『대수공간의 사상』의 보조정리
\ref{spaces-morphisms-lemma-smoothness-dimension-spaces}를 보라.
주석 \ref{stacks-geometry-remark-relative-dimension} (2)에 의해 여기서 사용한
$\dim_u (U_x)$의 정의가 그곳에서 사용한 정의와 일치함에 유의하자.
\end{proof}

\begin{lemma}
\label{stacks-geometry-lemma-dimension-for-stacks}
$\mathcal{X}$가 국소 뇌터 대수적 스택이고 $x \in |\mathcal{X}|$라 하자.
$U \to \mathcal{X}$를 대수공간에서 $\mathcal{X}$로 가는 매끄러운 사상이라
하고, $u$를 $x$로 사상되는 $|U|$의 임의의 점이라 하자. 그러면
$$
\dim_x(\mathcal{X}) =  \dim_u(U) - \dim_{u}(U_x)
$$
이다. 여기서 상대 차원 $\dim_u(U_x)$는 정의
\ref{stacks-geometry-definition-relative-dimension}에 의해 정의하고,
$x$에서 $\mathcal{X}$의 차원은 『스택의 성질』의 정의
\ref{stacks-properties-definition-dimension-at-point}에서와 같다.
\end{lemma}

\begin{proof}
보조정리 \ref{stacks-geometry-lemma-behaviour-of-dimensions-wrt-smooth-morphisms}를 사용하면
우변 $\dim_u(U) + \dim_u(U_x)$가 매끄러운 사상
$U \to \mathcal{X}$와 $u \in |U|$의 선택에 무관함을 확인할 수 있다.
세부사항은 생략한다. 특히 $U$가 스킴이라고 가정해도 좋다.
이 경우 $x$의 대표를 합성 $\Spec \kappa(u) \to U \to \mathcal{X}$로
택하여 $\dim_u(U_x)$를 계산할 수 있다. 여기서 첫 번째 사상은
$u \in U$를 상으로 하는 표준 사상이다.
$R = U \times_{\mathcal{X}} U$로 쓰고 $e : U \to R$을 대각사상이라 하자.
상대 차원이 밑변환 아래에서 불변이므로
$\dim_u(U_x) = \dim_{e(u)}(R_u)$이다. 따라서 우변은 원하는 대로
$\dim_u (U) - \dim_{e(u)}(R_u) = \dim_x(\mathcal{X})$와 같다.
\end{proof}

\begin{remark}
\label{stacks-geometry-remark-dimension-DM}
충분히 적절한(예를 들어 준분리인) Deligne--Mumford 스택에 대해서는
다시 $\dim_x(\mathcal{X})$가 위상적으로 정의한 양
$\dim_x |\mathcal{X}|$와 일치한다. 그러나 더 일반적인 Artin 스택에서는
보통 그렇지 않다. 예를 들어
$\mathcal{X} = [\mathbf{A}^1/\mathbf{G}_m]$라 하자
(어떤 체 위에서 $\mathbf{A}^1$에 대한 $\mathbf{G}_m$의 통상적인
곱셈 작용으로 몫을 취한다). 그러면 $|\mathcal{X}|$에는 두 점이 있고
하나는 다른 하나의 특수화이므로($\mathbf{A}^1$ 위의
$\mathbf{G}_m$의 두 궤도에 대응한다), 위상공간으로서 차원이 $1$이다.
하지만 두 점 $x \in |\mathcal{X}|$ 모두에 대해
$\dim_x (\mathcal{X}) = 0$이다.
(더 극단적인 예는 분류 공간 $[\Spec k/\mathbf{G}_m]$이며,
그 유일한 점에서의 차원은 $-1$이다.)
\end{remark}

\noindent
이제 정의 \ref{stacks-geometry-definition-relative-dimension}을 (국소 뇌터)
대수적 스택들 사이의 (국소적으로 유한형인) 사상으로 확장할 수 있다.

\begin{definition}
\label{stacks-geometry-definition-relative-dimension-for-stacks}
$f : \mathcal{T} \to \mathcal{X}$가 국소 뇌터 대수적 스택들 사이의
국소적으로 유한형인 사상이고, $t \in |\mathcal{T}|$가 상이
$x \in |\mathcal{X}|$인 점이면, $t$에서 $f$의 {\it 상대 차원}을
$\dim_t(\mathcal{T}_x)$로 표기하고 다음과 같이 정의한다.
$x$를 나타내는, 체의 스펙트럼을 원천으로 하는 사상
$\Spec k \to \mathcal{X}$를 택하고, $|\mathcal{T}|$로의 사영 아래에서
$t$로 사상되는 점
$t' \in |\mathcal{T} \times_{\mathcal{X}} \Spec k|$를 택한다
(그러한 점 $t'$는 『스택의 성질』의 보조정리
\ref{stacks-properties-lemma-points-cartesian}에 의해 존재한다. 그러면
$$
\dim_t(\mathcal{T}_x) = \dim_{t'}(\mathcal{T} \times_{\mathcal{X}} \Spec k ).
$$
\end{definition}

\noindent
$\mathcal{T}$와 $\mathcal{X}$가 대수적 스택이므로 섬유곱
$\mathcal{T}\times_{\mathcal{X}} \Spec k$는 대수적 스택이고,
『스택의 사상』의 보조정리
\ref{stacks-morphisms-lemma-locally-finite-type-locally-noetherian}에 의해
국소 뇌터이다. 따라서 이 정의에서 제안한 우변의 양은
『스택의 성질』의 정의
\ref{stacks-properties-definition-dimension-at-point}에 의해 정의된다.

\begin{remark}
\label{stacks-geometry-remark-dimension-tangent-space-well-defined}
표준적인 조작으로 $\dim_t(\mathcal{T}_x)$가 그 계산에 사용한 선택과
무관하게 잘 정의됨을 알 수 있다.
\end{remark}

\noindent
이제 상대 차원의 몇 가지 기본 성질을 확립한다. 이는 스킴의 사상의 경우에
대응하는 명제들을 명백히 일반화한 것이다.

\begin{lemma}
\label{stacks-geometry-lemma-base-change-invariance-of-relative-dimension}
국소 뇌터 스택들의 사상의 카르테시안 정사각형
$$
\xymatrix{
\mathcal{T}' \ar[d]\ar[r] & \mathcal{T} \ar[d] \\
\mathcal{X}' \ar[r] & \mathcal{X}
}
$$
이 주어졌고, 수직 사상들이 국소적으로 유한형이라고 하자.
$t' \in |\mathcal{T}'|$의 $|\mathcal{T}|$, $|\mathcal{X}'|$,
$|\mathcal{X}|$에서의 상을 각각 $t$, $x'$, $x$라 하면,
$\dim_{t'}(\mathcal{T}'_{x'}) = \dim_{t}(\mathcal{T}_x)$이다.
\end{lemma}

\begin{proof}
정의에 의해 양변을 같은 섬유곱의 차원으로 계산할 수 있다.
\end{proof}

\begin{lemma}
\label{stacks-geometry-lemma-behaviour-of-dimensions-wrt-smooth-morphisms-stacky}
$f: \mathcal{U} \to \mathcal{X}$가 국소 뇌터 대수적 스택들의
매끄러운 사상이고, $u \in |\mathcal{U}|$의 상이
$x \in |\mathcal{X}|$이면
$$
\dim_u (\mathcal{U}) = \dim_x(\mathcal{X}) + \dim_{u} (\mathcal{U}_x).
$$
\end{lemma}

\begin{proof}
원천이 스킴인 매끄러운 전사 $V \to \mathcal{U}$를 택하고,
$v\in |V|$를 $u$로 사상되는 점이라 하자.
합성 $V \to \mathcal{U} \to \mathcal{X}$도 매끄러우므로,
보조정리 \ref{stacks-geometry-lemma-behaviour-of-dimensions-wrt-smooth-morphisms}에 의해
$\dim_x(\mathcal{X}) = \dim_v(V) - \dim_v(V_x)$이고,
$\dim_u(\mathcal{U}) = \dim_v(V) - \dim_v(V_u)$이다. 따라서
$$
\dim_u(\mathcal{U}) - \dim_x(\mathcal{X}) = \dim_v (V_x) - \dim_v (V_u).
$$

\medskip\noindent
$x$의 대표 $\Spec k \to \mathcal{X}$를 택하고,
$v$ 위에 놓이는 점 $v' \in | V \times_{\mathcal{X}} \Spec k|$를 택하자.
$|\mathcal{U}\times_{\mathcal{X}} \Spec k|$에서 그 상을 $u'$라 하자.
그러면 정의에 의해
$\dim_u(\mathcal{U}_x) = \dim_{u'}(\mathcal{U}\times_{\mathcal{X}} \Spec k)$이고,
$\dim_v(V_x) = \dim_{v'}(V\times_{\mathcal{X}} \Spec k)$이다.

\medskip\noindent
이제 $V\times_{\mathcal{X}} \Spec k \to
\mathcal{U}\times_{\mathcal{X}}\Spec k$는 그러한 사상의 밑변환이므로
매끄러운 전사이고, 그 원천은 대수공간이다
($V$와 $\Spec k$가 스킴이고 $\mathcal{X}$가 대수적 스택이기 때문이다).
따라서 다시 정의에 의해
\begin{align*}
\dim_{u'}(\mathcal{U}\times_{\mathcal{X}} \Spec k)
& =
\dim_{v'}(V\times_{\mathcal{X}} \Spec k) -
\dim_{v'}(V \times_{\mathcal{X}} \Spec k)_{u'}) \\
& = \dim_v(V_x) -
\dim_{v'}( (V\times_{\mathcal{X}} \Spec k)_{u'}).
\end{align*}
이제 $V\times_{\mathcal{X}} \Spec k \cong
V\times_{\mathcal{U}} (\mathcal{U}\times_{\mathcal{X}} \Spec k)$이므로,
보조정리 \ref{stacks-geometry-lemma-base-change-invariance-of-relative-dimension}에 의해
$\dim_{v'}((V\times_{\mathcal{X}} \Spec k)_{u'})  = \dim_v(V_u)$이다.
모든 것을 합치면 요구한 대로
$$
\dim_u(\mathcal{U}) - \dim_x(\mathcal{X}) =
\dim_u(\mathcal{U}_x),
$$
를 얻는다.
\end{proof}

\begin{lemma}
\label{stacks-geometry-lemma-relative-dimension-is-semi-continuous}
$f: \mathcal{T} \to \mathcal{X}$를 대수적 스택들의
국소적으로 유한형인 사상이라 하자.
\begin{enumerate}
\item
함수 $t \mapsto \dim_t(\mathcal{T}_{f(t)})$는
$|\mathcal{T}|$ 위에서 상반연속이다.
\item $f$가 매끄러우면 함수
$t \mapsto \dim_t(\mathcal{T}_{f(t)})$는
$|\mathcal{T}|$ 위에서 국소 상수이다.
\end{enumerate}
\end{lemma}

\begin{proof}
먼저 $\mathcal{T}$가 스킴 $T$라고 하자.
원천이 스킴인 매끄러운 전사 $U \to \mathcal{X}$를 택하고,
$T' = T \times_{\mathcal{X}} U$로 두자.
$f': T' \to U$를 $U$ 위로 당겨온 $f$라 하고,
$g: T' \to T$를 사영이라 하자.

\medskip\noindent
보조정리 \ref{stacks-geometry-lemma-base-change-invariance-of-relative-dimension}에 의해
$t' \in T'$에 대해
$\dim_{t'}(T'_{f'(t')}) = \dim_{g(t')}(T_{f(g(t'))})$이다.
한편 $g$는 매끄러운 전사의 밑변환이므로 매끄럽고 전사이며,
$g$가 바탕 위상공간들 위에 유도하는 사상은
『대수공간의 성질』의 보조정리
\ref{spaces-properties-lemma-topology-points}에 의해 연속이고 열려 있으며
전사이다. 따라서 (1)은 사상 $f'$에 대해 『대수공간의 사상』의 보조정리
\ref{spaces-morphisms-lemma-openness-bounded-dimension-fibres}에서 따르고,
(2)는 『스킴의 사상』의 보조정리
\ref{morphisms-lemma-flat-finite-presentation-CM-fibres-relative-dimension} 또는
『스킴의 사상』의 보조정리
\ref{morphisms-lemma-smooth-omega-finite-locally-free}에서 따른다는 점만
확인하면 충분하다(두 결과 모두 스킴에 대한 결과를 주며, 그 결과에서
대수공간에 대한 유사한 결과는 『대수공간의 사상』의 보조정리
\ref{spaces-morphisms-lemma-openness-bounded-dimension-fibres}에서와
똑같이 추론할 수 있다.

\medskip\noindent
이제 일반적인 경우로 돌아가 원천이 스킴인 매끄러운 전사
$h:V \to \mathcal{T}$를 택하자. $v \in V$이면 사실상 정의에 의해
$$
\dim_{h(v)}(\mathcal{T}_{f(h(v))}) =
\dim_{v}(V_{f(h(v))}) - \dim_{v}(V_{h(v)}).
$$
$V$가 스킴이므로, 이 등식 우변의 첫 번째 항이 상반연속이고
($f$가 매끄러우면 심지어 국소 상수이고), 두 번째 항은 실제로
국소 상수임을 증명했다. 따라서 그 차는 상반연속이고
($f$가 매끄러우면 국소 상수이고), 그러므로 함수
$\dim_{h(v)}(\mathcal{T}_{f(h(v))})$는 $|V|$ 위에서 상반연속이며
($f$가 매끄러우면 국소 상수이다).
사상 $|V| \to |\mathcal{T}|$가 열리고 전사이므로 보조정리가 따른다.
\end{proof}

\noindent
논의를 계속하기 전에 스킴의 차원론과 관련된 두 보조정리를 증명한다.

\medskip\noindent
첫 번째 보조정리의 맥락을 설명하자. $X$가 유한 차원 스킴이면
$\dim X$는 차원 $\dim_x X$들의 상한으로 정의되므로
$\dim_x X = \dim X$인 점 $x \in X$가 존재한다.
다음 보조정리는 더 나아가 $x$를 유한형 점으로 택할 수 있음을 보인다.

\begin{lemma}
\label{stacks-geometry-lemma-dimension-achieved-by-finite-type-point}
$X$가 유한 차원 스킴이면 $\dim_x X = \dim X$인
닫힌(따라서 유한형인) 점 $x \in X$가 존재한다.
\end{lemma}

\begin{proof}
$d = \dim X$로 두고, $X$의 기약 닫힌부분집합들의 극대 진감소 사슬
\begin{equation}
\label{stacks-geometry-equation-maximal-chain}
Z_0 \supset Z_1 \supset \ldots \supset Z_d.
\end{equation}
을 택하자. 부분집합 $Z_d$는 $X$의 극소 기약 닫힌부분집합이므로,
$Z_d$의 모든 점은 $Z_d$의 일반점이다. 스킴 $X$의 바탕 위상공간은
소버 공간이므로, $Z_d$는 하나의 닫힌점 $x \in X$로 이루어진
한원소 집합이라고 결론 내린다. $U$가 $x$의 임의의 근방이면 사슬
$$
U\cap Z_0 \supset U\cap Z_1 \supset \ldots \supset U\cap Z_d = Z_d =
\{x\}
$$
은 $U$의 기약 닫힌부분집합들의 진감소 사슬이고, 따라서
$\dim U \geq d$이다. 그러므로 $\dim_x X \geq d$를 얻는다.
반대 부등식은 명백하므로 보조정리가 증명되었다.
\end{proof}

\noindent
다음 보조정리는 몇 가지 온화한 추가 가정을 만족하는 기약 스킴에서
$\dim_x X$가 {\it 상수} 함수임을 보인다.

\begin{lemma}
\label{stacks-geometry-lemma-constancy-of-dimension}
$X$가 유한 차원의 기약이고 제이컵슨이며 코테너리인 국소 뇌터 스킴이면,
$X$의 공집합이 아닌 모든 열린부분집합 $U$에 대해
$\dim U = \dim X$이다. 동치로 $\dim_x X$는 $X$ 위의 상수 함수이다.
\end{lemma}

\begin{proof}
두 주장의 동치는 정의에서 곧바로 따른다.
그러므로 $U\subset X$가 공집합이 아닌 열린부분집합이라고 하자.
분명히 $\dim U \leq \dim X$이며, $\dim U \geq \dim X$를 보여야 한다.
$d = \dim X$로 쓰고, $X$의 기약 닫힌부분집합들의 극대 진감소 사슬
$$
X = Z_0 \supset Z_1 \supset \ldots \supset Z_d.
$$
를 택하자. $X$가 제이컵슨이므로 극소 기약 닫힌부분집합 $Z_d$는
어떤 닫힌점 $x$에 대해 $\{x\}$와 같다.

\medskip\noindent
$x \in U$이면
$$
U = U \cap Z_0  \supset U\cap Z_1 \supset \ldots \supset
U\cap Z_d = \{x\}
$$
은 $U$의 기약 닫힌부분집합들의 진감소 사슬이다.
따라서 요구한 대로 $\dim U \geq d$이다.
그러므로 이제 $x \not\in U$라고 가정해도 좋다.

\medskip\noindent
평탄 사상 $\Spec \mathcal{O}_{X,x} \to X$를 생각하자.
$X$의 공집합이 아닌(따라서 조밀한) 열린부분집합 $U$는
열린부분집합 $V \subset \Spec \mathcal{O}_{X,x}$로 당겨진다.
$U$를 공집합이 아닌 준콤팩트, 따라서 뇌터인 열린부분집합으로 대체하여,
포함사상 $U \to X$가 준콤팩트 사상이라고 가정해도 좋다.
준콤팩트 사상의 스킴론적 상의 형성은 평탄 밑변환과 가환하므로
『스킴의 사상』의 보조정리
\ref{morphisms-lemma-flat-base-change-scheme-theoretic-image}에 의해
$V$는 $\Spec \mathcal{O}_{X,x}$에서 조밀하고, 특히 공집합이 아니며,
물론 $x \not\in V$이다. (여기서는
$\Spec \mathcal{O}_{X,x}$의 닫힌점도 $x$로 나타낸다. 그 상이 주어진 점
$x \in X$와 같기 때문이다.)
이제 『스킴의 성질』의 보조정리
\ref{properties-lemma-complement-closed-point-Jacobson}에 의해
$\Spec \mathcal{O}_{X,x} \setminus \{x\}$는 제이컵슨이고,
따라서 $V$는 $\Spec \mathcal{O}_{X,x} \setminus \{x\}$의
닫힌점 $z$를 포함한다. $X$ 안에서 $z$의 상의 폐포는
$x$를 포함하고 $U$와의 교집합이 공집합이 아닌 $X$의
기약 닫힌부분집합 $Z$이며, $Z$에 진포함되고 $\{x\}$를 진포함하는
기약 닫힌부분집합은 없다
($\Spec \mathcal{O}_{X,x}$로의 당김이 $x$를 포함하는 $X$의
기약 닫힌부분집합과 $\Spec \mathcal{O}_{X,x}$의 기약 닫힌부분집합
사이의 전단사를 유도하기 때문이다).
$U \cap Z$는 공집합이 아닌 $U$의 닫힌부분집합이므로,
$X$가 제이컵슨이라는 사실에 의해 $X$에서 닫힌 점 $u$를 포함한다.
또한 $U\cap Z$는 기약집합 $Z$의 공집합이 아닌(따라서 조밀한)
열린부분집합이고, 이 기약집합은 $U$에 속하지 않는 점 $x$를 포함하므로
포함관계 $\{u\} \subset U\cap Z$는 진포함이다.

\medskip\noindent
$X$가 코테너리이므로 사슬
$$
X = Z_0 \supset Z \supset \{x\} = Z_d
$$
은 길이 $d+1$인 사슬로 세분할 수 있고, 이는 다음 형태여야 한다.
$$
X = Z_0 \supset W_1 \supset \ldots \supset W_{d-1} = Z \supset \{x\} = Z_d.
$$
$U\cap Z$가 공집합이 아니므로
$$
U = U \cap Z_0 \supset U \cap W_1\supset \ldots \supset U\cap W_{d-1}
= U\cap Z \supset \{u\}
$$
는 $U$의 기약 닫힌부분집합들의 길이 $d+1$인 진감소 사슬이다.
따라서 요구한 대로 $\dim U \geq d$이다.
\end{proof}

\noindent
보조정리 \ref{stacks-geometry-lemma-constancy-of-dimension}의 스택론적 유사물을
아래의 보조정리 \ref{stacks-geometry-lemma-irreducible-implies-equidimensional}에서
증명할 것이다. 그러나 그에 앞서 추가 정의를 도입해야 한다.
스킴이 코테너리라는 개념은 에탈 국소적이지 않기 때문에 이것이 필요하다
(『대수학』의 주석
\ref{algebra-remark-universally-catenary-does-not-descend}의 예를 보라.
이 때문에 대수공간이나 대수적 스택이 코테너리라는 말의 의미를
정의하기가 어렵다(\cite[page 3]{Osserman}의 논의를 보라).
차원론의 몇몇 측면에서는 다음 정의가 코테너리 대수적 스택이라는
빠진 개념을 잘 대신하는 듯하다.

\begin{definition}
\label{stacks-geometry-definition-pseudo-catenary}
국소 뇌터 대수적 스택 $\mathcal{X}$에 대해, 원천이 보편적으로
코테너리인 스킴인 매끄러운 전사 $U \to \mathcal{X}$가 존재하면
이를 {\it 유사 코테너리}라고 한다.
\end{definition}

\begin{example}
\label{stacks-geometry-example-pseudo-catenary}
$\mathcal{X}$가 보편적으로 코테너리인 국소 뇌터 스킴 $S$ 위에서
국소적으로 유한형이고, $U\to \mathcal{X}$가 원천이 스킴인
매끄러운 전사라고 하자. 그러면 합성 $U \to \mathcal{X} \to S$는
국소적으로 유한형이므로 『스킴의 사상』의 보조정리
\ref{morphisms-lemma-universally-catenary-local}에 의해 $U$는
보편적으로 코테너리이다. 따라서 $\mathcal{X}$는 유사 코테너리이다.
\end{example}

\noindent
다음 보조정리는 유사 코테너리라는 성질이 유한형 사상을 따라
전달됨을 보인다.

\begin{lemma}
\label{stacks-geometry-lemma-catenary-covers}
$\mathcal{X}$가 유사 코테너리인 국소 뇌터 대수적 스택이고,
$\mathcal{Y} \to \mathcal{X}$가 국소적으로 유한형인 사상이면,
원천이 보편적으로 코테너리인 스킴인 매끄러운 전사
$V \to \mathcal{Y}$가 존재한다. 따라서 $\mathcal{Y}$도
유사 코테너리이다.
\end{lemma}

\begin{proof}
가정에 의해 원천이 보편적으로 코테너리인 스킴인 매끄러운 전사
$U \to \mathcal{X}$를 찾을 수 있다. 밑변환
$U\times_{\mathcal{X}} \mathcal{Y}$는 대수적 스택이다.
$V \to U\times_{\mathcal{X}} \mathcal{Y}$를 원천이 스킴인
매끄러운 전사라 하자. 그러면 합성
$V \to U\times_{\mathcal{X}} \mathcal{Y} \to \mathcal{Y}$는
매끄러운 전사들의 합성이므로 매끄럽고 전사이다. 한편 사상
$V \to U\times_{\mathcal{X}} \mathcal{Y} \to U$는
국소적으로 유한형인 사상들의 합성이므로 국소적으로 유한형이다.
$U$가 보편적으로 코테너리이므로 『스킴의 사상』의 보조정리
\ref{morphisms-lemma-universally-catenary-local}에 의해
주장한 대로 $V$도 보편적으로 코테너리이다.
\end{proof}

\noindent
이제 (어떤 국소 뇌터 스택 $\mathcal{X}$에 대해) $|\mathcal{X}|$ 위의
함수 $\dim_x(\mathcal{X})$가 $|\mathcal{X}|$의 기약 성분들에 대해
어떻게 행동하는지, 그리고 이와 관련된 여러 주제를 연구한다.

\begin{lemma}
\label{stacks-geometry-lemma-irreducible-implies-equidimensional}
$\mathcal{X}$가 제이컵슨이고 유사 코테너리인 국소 뇌터 대수적 스택이며
$|\mathcal{X}|$가 기약이면, $\dim_x(\mathcal{X})$는
$|\mathcal{X}|$ 위의 상수 함수이다.
\end{lemma}

\begin{proof}
$\dim_x(\mathcal{X})$가 $|\mathcal{X}|$ 위에서 국소 상수임을 보이면 충분하다.
그러면 $|\mathcal{X}|$가 기약이어서 연결이므로 반드시 상수이기 때문이다.
$\mathcal{X}$가 유사 코테너리이므로, $U$가 보편적으로 코테너리인 스킴인
매끄러운 전사 $U \to \mathcal{X}$를 찾을 수 있다.
$\{U_i\}$가 준콤팩트 열린부분스킴들로 이루어진 $U$의 덮개이면
$U$를 $\coprod U_i,$,로 대체할 수 있으며, 함수
$u \mapsto \dim_{f(u)}(\mathcal{X})$가 $U_i$ 위에서 국소 상수임을
보이면 충분하다. 이를 한 번에 하나의 $U_i$에 대해 확인하므로
이제 아래첨자를 생략하고 $U_i$ 대신 단순히 $U$라고 쓴다.
$U$는 준콤팩트이므로 유한 개의 기약 성분, 이를테면
$T_1 \cup \ldots \cup T_n$의 합집합이다.
각 $T_i$는 제이컵슨이고 코테너리인 국소 뇌터 스킴 $U$의
닫힌부분스킴이므로, 그 자신도 제이컵슨이고 코테너리이며 국소 뇌터이다.

\medskip\noindent
보조정리 \ref{stacks-geometry-lemma-behaviour-of-dimensions-wrt-smooth-morphisms}에 의해
$\dim_{f(u)}(\mathcal{X}) = \dim_{u}(U) - \dim_{u}(U_{f(u)})$이다.
$f$가 매끄러우므로 보조정리
\ref{stacks-geometry-lemma-relative-dimension-is-semi-continuous} (2)에 의해
우변의 두 번째 항은 $U$ 위에서 국소 상수이다. 따라서
$\dim_u(U)$가 $U$ 위에서 국소 상수임을 보여야 한다.
$\dim_u(U)$는 $u$를 포함하는 $U$의 성분들 $T_i$에 대한
$\dim_u T_i$의 최댓값이므로, 점 $u$가 서로 다른 두 성분
$T_i$와 $T_j$($i \neq j$)에 놓이면 $\dim_u T_i = \dim_u T_j$임을
보이는 것으로 충분하다. 그런 다음 기약이고 제이컵슨이며 코테너리인
국소 뇌터 스킴 $T$ 위에서 $t\mapsto \dim_t T$가 상수 함수라는 점을
사용한다(보조정리 \ref{stacks-geometry-lemma-constancy-of-dimension}에서 따른다).

\medskip\noindent
$V = T_i \setminus (\bigcup_{i' \neq i} T_{i'})$ 및
$W = T_j \setminus (\bigcup_{i' \neq j} T_{i'})$로 두자.
그러면 $V$와 $W$는 각각 $U$의 공집합이 아닌 열린부분집합이고,
따라서 $|\mathcal{X}|$에서 각각의 상은 공집합이 아닌 열린집합이다.
$|\mathcal{X}|$가 기약이므로 $|\mathcal{X}|$의 이 두 열린부분집합의 교집합은
공집합이 아니다. 이 교집합에 놓이는 점을 $x$라 하고,
$x$로 사상되는 점들을 $v \in V$, $w\in W$라 하자. 그러면
$$
\dim T_i = \dim V = \dim_v (U) = \dim_x (\mathcal{X}) + \dim_v (U_x)
$$
이고 마찬가지로
$$
\dim T_j = \dim W = \dim_w (U) = \dim_x (\mathcal{X}) + \dim_w (U_x).
$$
함수 $u \mapsto \dim_u (U_{f(u)})$는 $U$ 위에서 국소 상수이고,
$T_i \cup T_j$는 교집합이 공집합이 아닌 두 기약, 따라서 연결인 집합의
합집합이므로 연결이다. 따라서 $\dim_v (U_x) = \dim_w(U_x)$이고,
앞의 두 등식을 비교하면 요구한 대로 $\dim T_i = \dim T_j$이다.
\end{proof}

\begin{lemma}
\label{stacks-geometry-lemma-closed-immersions}
$\mathcal{Z} \hookrightarrow \mathcal{X}$가 국소 뇌터 대수적 스택들의
닫힌 몰입이고, $z \in |\mathcal{Z}|$의 상이
$x \in |\mathcal{X}|$이면
$\dim_z (\mathcal{Z}) \leq \dim_x(\mathcal{X})$이다.
\end{lemma}

\begin{proof}
원천이 스킴인 매끄러운 전사 $U\to \mathcal{X}$를 택한다.
밑변환된 사상 $V = U\times_{\mathcal{X}} \mathcal{Z} \to \mathcal{Z}$도
매끄럽고 전사이며, 사영 $V \to U$는 닫힌 몰입이다.
$v \in |V|$가 $z \in |\mathcal{Z}|$로 사상된다고 하고,
$u$를 $|U|$에서 $v$의 상이라 하자. 그러면 분명히
$\dim_v(V) \leq \dim_u(U)$이고, 보조정리
\ref{stacks-geometry-lemma-base-change-invariance-of-relative-dimension}에 의해
$\dim_v (V_z) = \dim_u(U_x)$이다. 따라서 주장한 대로
$$
\dim_z(\mathcal{Z})  = \dim_v(V) - \dim_v(V_z)
\leq \dim_u(U) - \dim_u(U_x) = \dim_x(\mathcal{X}),
$$
이다.
\end{proof}
\begin{lemma}
\label{stacks-geometry-lemma-dimension-via-components}
$\mathcal{X}$가 국소 뇌터 대수적 스택이고 $x \in |\mathcal{X}|$이면
$\dim_x(\mathcal{X}) = \sup_{\mathcal{T}} \{ \dim_x(\mathcal{T}) \} $,
여기서 $\mathcal{T}$는 $x$를 지나는 $|\mathcal{X}|$의 모든
기약 성분을 그 유도된 축소 구조와 함께 달린다.
\end{lemma}

\begin{proof}
보조정리 \ref{stacks-geometry-lemma-closed-immersions}에 의해 점 $x$를 지나는 각
기약 성분 $\mathcal{T}$에 대해
$\dim_x (\mathcal{T}) \leq \dim_x(\mathcal{X})$이다.
따라서 보조정리를 증명하려면
\begin{equation}
\label{stacks-geometry-equation-desired-inequality}
\dim_x(\mathcal{X}) \leq
\sup_{\mathcal{T}} \{\dim_x(\mathcal{T})\}.
\end{equation}
임을 보이면 충분하다. $U\to\mathcal{X}$를 스킴에 의한 매끄러운
덮개라 하자. $T$가 $U$의 기약 성분이면 $\mathcal{T}$를
$\mathcal{X}$ 안에서 그 상의 폐포라 하자. 이는 $\mathcal{X}$의
기약 성분이다. $u \in U$를 $x$로 사상되는 점이라 하자. 그러면
$\dim_x(\mathcal{X})=\dim_uU-\dim_uU_x=\sup_T\dim_uT-\dim_uU_x$이다.
여기서 상한은 $u$를 지나는 $U$의 기약 성분들에 걸쳐 취한다.
상한이 달성되는 성분 $T$를 택하고
$\dim_x(\mathcal{T})=\dim_uT-\dim_u T_x$임에 유의하자.
이제 원하는 부등식 (\ref{stacks-geometry-equation-desired-inequality})은 명백한 부등식
$\dim_u T_x \leq \dim_u U_x$에서 따른다.
($\Spec k \to \mathcal{X}$가 $x$의 대표이면
$T\times_{\mathcal{X}} \Spec k$는
$U\times_{\mathcal{X}} \Spec k$의 닫힌 부분공간임에 유의하자.)
\end{proof}

\begin{lemma}
\label{stacks-geometry-lemma-dimension-at-finite-type-point}
$\mathcal{X}$가 국소 뇌터 대수적 스택이고 $x \in |\mathcal{X}|$이면,
$x$를 포함하는 $\mathcal{X}$의 임의의 열린 부분스택 $\mathcal{V}$에 대해
$\dim_{x_0}(\mathcal{X}) = \dim_x(\mathcal{V})$인 유한형 점
$x_0 \in |\mathcal{V}|$가 존재한다.
\end{lemma}

\begin{proof}
원천이 스킴인 매끄러운 전사 $f:U \to \mathcal{X}$를 택하고,
함수 $u \mapsto \dim_{f(u)}(\mathcal{X})$를 생각하자.
$f$가 매끄러우므로 $f$가 유도하는 사상
$|U| \to |\mathcal{X}|$는 열리고, 가정에 의해 전사이며,
$|\mathcal{X}|$의 유한형 점의 정의 자체에 의해 유한형 점을
유한형 점으로 보낸다. 따라서 임의의 $u \in U$와 $u$의 임의의
열린 근방에 대해, 이 근방 안에
$\dim_{f(u_0)}(\mathcal{X}) = \dim_{f(u)}(\mathcal{X})$인 유한형 점
$u_0$가 존재함을 보이면 충분하다.
문제를 이렇게 다시 표현하면 $f$의 전사성은 더 이상 필요하지 않으므로,
$U$를 문제의 점 $u$의 열린 근방으로 대체할 수 있다.
따라서 각 $u \in U$에 대해
$\dim_{f(u_0)}(\mathcal{X}) = \dim_{f(u)}(\mathcal{X})$인 유한형 점
$u_0 \in U$가 존재함을 보이는 문제로 환원된다.
보조정리 \ref{stacks-geometry-lemma-behaviour-of-dimensions-wrt-smooth-morphisms}에 의해
$\dim_{f(u)}(\mathcal{X}) = \dim_u(U) - \dim_u(U_{f(u)})$이고,
$\dim_{f(u_0)}(\mathcal{X}) = \dim_{u_0}(U) - \dim_{u_0}(U_{f(u_0)})$이다.
$f$가 매끄러우므로 $u_0$가 $U$ 위에서 변할 때 식
$\dim_{u_0}(U_{f(u_0)})$은 국소 상수이다
(보조정리 \ref{stacks-geometry-lemma-relative-dimension-is-semi-continuous} (2)).
그러므로 필요하면 $u$ 주위로 $U$를 더 축소하여 이 식이 상수라고
가정해도 좋다. 따라서 $\dim_{u_0}(U) = \dim_u(U)$인
유한형 점 $u_0 \in U$를 찾을 수 있음을 보이는 문제가 된다.
정의에 의해 $\dim_u U$는 $U$ 안에서 $u$의 열린 근방 $V$들에 대한
$\dim V$의 최솟값이므로, $u$ 주위로 $U$를 더 축소하여
$\dim_u U = \dim U$라고 해도 좋다. 이제 원하는 점 $u_0$의 존재는
보조정리 \ref{stacks-geometry-lemma-dimension-achieved-by-finite-type-point}에서 따른다.
\end{proof}

\begin{lemma}
\label{stacks-geometry-lemma-monomorphing-a-component-in-of-the-right-dimension}
$\mathcal{T} \hookrightarrow \mathcal{X}$를 대수적 스택들의
국소적으로 유한형인 단사상이라 하자. $\mathcal{X}$와 (따라서)
$\mathcal{T}$는 제이컵슨이고 유사 코테너리이며 국소 뇌터라고 하자.
더 나아가 $\mathcal{T}$는 어떤 (유한한) 차원 $d$의 기약 스택이고,
$\mathcal{X}$는 축소되며 차원이 $d$ 이하라고 하자.
그러면 $\mathcal{T}$의 공집합이 아닌 열린 부분스택 $\mathcal{V}$가
존재하여, 유도된 단사상 $\mathcal{V} \hookrightarrow \mathcal{X}$는
$\mathcal{V}$를 $\mathcal{X}$의 한 기약 성분의 열린부분집합과
동일시하는 열린 몰입이다.
\end{lemma}

\begin{proof}
원천이 스킴인 매끄러운 전사 $f:U \to \mathcal{X}$를 택한다.
$\mathcal{X}$가 축소되었으므로 이 스킴도 반드시 축소되어 있다.
$U' = \mathcal{T}\times_{\mathcal{X}} U$로 쓰자.
밑변환된 사상 $U' \to U$는 대수공간들의 국소적으로 유한형인 단사상이고,
따라서 『대수공간의 사상』의 보조정리
\ref{spaces-morphisms-lemma-locally-quasi-finite-separated-representable}와
\ref{spaces-morphisms-lemma-monomorphism-loc-finite-type-loc-quasi-finite}에 의해
표현가능하다. $U$가 스킴이므로 $U'$도 스킴이다.
사영 $f': U' \to \mathcal{T}$도 매끄러운 전사이다.
$u' \in U'$라 하고 그 상을 $u \in U$라 하자.
보조정리 \ref{stacks-geometry-lemma-base-change-invariance-of-relative-dimension}에 의해
$\dim_{u'}(U'_{f(u')}) = \dim_u(U_{f(u)})$이고,
보조정리 \ref{stacks-geometry-lemma-irreducible-implies-equidimensional} 및
$\mathcal{T}$와 $\mathcal{X}$에 대한 우리의 가정에 의해
$\dim_{f'(u')}(\mathcal{T}) =d
\geq \dim_{f(u)}(\mathcal{X})$이다. 따라서
\begin{equation}
\label{stacks-geometry-equation-dim-inequality}
\dim_{u'} (U') = \dim_{u'} (U'_{f(u')}) + \dim_{f'(u')}(\mathcal{T})
\\
\geq \dim_u (U_{f(u)}) + \dim_{f(u)}(\mathcal{X}) = \dim_u (U).
\end{equation}
$U' \to U$는 국소적으로 유한형인 단사상이므로 특히 비분기이다.
따라서 비분기 사상의 에탈 국소 구조에 관한 『에탈 사상』의 보조정리
\ref{etale-lemma-finite-unramified-etale-local}에 의해 가환 도식
$$
\xymatrix{
V' \ar[r]\ar[d] & V \ar[d] \\
U' \ar[r] & U
}
$$
을 찾을 수 있다. 여기서 스킴 $V'$는 공집합이 아니고,
수직 화살표들은 에탈이며, 위쪽 수평 화살표는 닫힌 몰입이다.
$V$를 그 상이 $U'$의 상과 공집합이 아닌 교집합을 갖는
준콤팩트 열린부분집합으로 대체하고 $V'$을 $V$의 역상으로 대체하면,
$V$(따라서 $V'$도)가 준콤팩트이라고 더 가정할 수 있다.
$V$는 국소 뇌터이기도 하므로 뇌터이고, 따라서 유한 개의
기약 성분의 합집합이다.

\medskip\noindent
에탈 사상은 점별 차원을 보존하므로
『내림』의 보조정리 \ref{descent-lemma-dimension-at-point-local}와
(\ref{stacks-geometry-equation-dim-inequality})에 의해, $v \in V$로 사상되는
임의의 점 $v' \in V'$에 대해
$\dim_{v'}( V') \geq \dim_v(V)$이다.
특히 $V'$의 상은 $V$의 서로 다른 두 기약 성분의 교집합에
포함될 수 없다. 따라서 $V'$과 공집합이 아닌 교집합을 갖는
$V$의 기약 열린부분집합을 적어도 하나 찾을 수 있다.
$V$를 이 부분집합으로 대체하여 $V$가 정수적이라고 가정해도 좋다
($V$는 축소되고 기약이기 때문이다).
앞의 차원 부등식에서 닫힌 몰입 $V' \hookrightarrow V$가 실제로
동형사상이라고 결론 내린다.
$W$를 $U'$ 안에서 $V'$의 상이라 하자. 그러면 에탈 사상은 열리므로
$W$는 $U'$의 공집합이 아닌 열린부분집합이다. 유도된 단사상
$W \to U$는 원천 위에서 에탈 국소적으로, 즉 $V'$으로 당긴 뒤 에탈이므로
에탈이고, 따라서 에탈 단사상으로서 열린 몰입이다.
그러므로 $\mathcal{V}$를 $\mathcal{T}$ 안에서 $W$의 상이라 하면,
$\mathcal{V}$는 $\mathcal{T}$의 조밀한(동치로 공집합이 아닌)
열린 부분스택이며, 그 상은 $\mathcal{X}$의 한 기약 성분에서 조밀하다.
마지막으로 사상 $\mathcal{V} \to \mathcal{X}$가 매끄럽다는 점을
언급한다(매끄러운 사상 $W\to \mathcal{V}$와의 합성이 매끄럽기 때문이다).
이 사상은 단사상이기도 하므로 열린 몰입이다.
\end{proof}
\begin{lemma}
\label{stacks-geometry-lemma-dims-of-images}
$f: \mathcal{T} \to \mathcal{X}$를 제이컵슨이고 유사 코테너리이며
국소 뇌터인 대수적 스택들의 국소적으로 유한형인 사상이라 하자.
원천은 기약이고 표적은 준분리라고 하며,
$\mathcal{Z} \hookrightarrow \mathcal{X}$를 $\mathcal{T}$의
스킴론적 상이라 하자. 그러면 모든 $t \in |T|$에 대해
$\dim_t( \mathcal{T}_{f(t)}) \geq \dim \mathcal{T}  - \dim \mathcal{Z}$이고,
그 위에서 등식이 성립하는 $|\mathcal{T}|$의 공집합이 아닌
(동치로 조밀한) 열린부분집합이 존재한다.
\end{lemma}

\begin{proof}
$\mathcal{X}$를 $\mathcal{Z}$로 대체하여 $f$가 스킴론적으로 우세하고
$\mathcal{X}$가 기약이라고 가정할 수 있고, 그렇게 한다.
섬유 차원의 상반연속성(보조정리
\ref{stacks-geometry-lemma-relative-dimension-is-semi-continuous} (1))에 의해,
$\mathcal{T}$의 어떤 공집합이 아닌 열린 부분스택에 속하는 $t$에 대해
등식 $\dim_t( \mathcal{T}_{f(t)}) =\dim \mathcal{T}  - \dim \mathcal{Z}$이
성립함을 증명하면 충분하다.
이 때문에 논증에서 언제나 $\mathcal{T}$를 공집합이 아닌
열린 부분스택으로 자유롭게 대체할 수 있다.

\medskip\noindent
$T' \to \mathcal{T}$를 원천이 스킴인 매끄러운 전사라 하고,
$T$를 $T'$의 공집합이 아닌 준콤팩트 열린부분집합이라 하자.
$\mathcal{Y}$가 준분리이므로 $T \to \mathcal{Y}$는 준콤팩트이다
(『스택의 사상』의 보조정리
\ref{stacks-morphisms-lemma-quasi-compact-permanence}를 사상
$T \to \mathcal{Y} \to \Spec \mathbf{Z}$에 적용하라).
따라서 $\mathcal{T}$를 $\mathcal{T}$ 안에서 $T$의 상으로 대체하면,
사상 $f:\mathcal{T} \to \mathcal{X}$가 (『스택의 사상』의 보조정리
\ref{stacks-morphisms-lemma-surjection-from-quasi-compact}에 호소하여
준콤팩트라고 가정할 수 있다.

\medskip\noindent
$U$가 스킴인 매끄러운 전사 $U \to \mathcal{X}$를 택하면,
보조정리 \ref{stacks-geometry-lemma-map-of-components}에 의해 $U$의 기약 열린부분집합
$V$로서 $V \to \mathcal{X}$가 매끄럽고 스킴론적으로 우세한 것을
찾을 수 있다. 준콤팩트 사상의 스킴론적 우세성은 평탄 밑변환으로
보존되므로, 스킴론적으로 우세한 사상 $f$의 밑변환
$\mathcal{T} \times_{\mathcal{X}} V \to V$도 스킴론적으로 우세하다.
이 섬유곱으로의 매끄러운 전사를 갖는 스킴을 $Z$라 하자.
그러면 $Z \to \mathcal{T} \times_{\mathcal{X}} V \to V$도
스킴론적으로 우세하다. 따라서 $V$를 스킴론적으로 지배하는
$Z$의 기약 성분 $C$를 찾을 수 있다.
합성 $Z \to \mathcal{T}\times_{\mathcal{X}} V \to \mathcal{T}$는
매끄럽고 $\mathcal{T}$는 기약이므로, 보조정리
\ref{stacks-geometry-lemma-map-of-components}에 의해 원천의 모든 기약 성분은
$|\mathcal{T}|$에서 조밀한 상을 갖는다.
이제 $C$를 $Z$의 다른 모든 기약 성분과 서로소인 공집합이 아닌
열린부분집합 $W$로 대체하고, 이어서 $\mathcal{T}$와 $\mathcal{X}$를
각각 $W$와 $V$의 상으로 대체한다
(보조정리 \ref{stacks-geometry-lemma-irreducible-implies-equidimensional}을 적용하면
이것이 $\mathcal{T}$나 $\mathcal{X}$의 차원을 바꾸지 않음을 알 수 있다).
$\mathcal{W}$를 사상
$W \to \mathcal{T}\times_{\mathcal{X}} V$의 상이라 하자.
그러면 $W \to \mathcal{T}\times_{\mathcal{X}} V$가 매끄러우므로
$\mathcal{W}$는 $\mathcal{T}\times_{\mathcal{X}} V$에서 열려 있고,
기약 스킴의 상이므로 기약이다. 이로써 가환 도식
$$
\xymatrix{
W \ar[dr] \ar[r]  & \mathcal{W} \ar[r] \ar[d] & V \ar[d] \\
& \mathcal{T} \ar[r] & \mathcal{X}
}
$$
을 얻는다. 여기서 $W$와 $V$는 스킴이고, 수직 화살표들은
매끄럽고 전사이며, 대각 화살표들과 왼쪽 위 수평 화살표는 매끄럽고,
유도된 사상 $\mathcal{W} \to \mathcal{T}\times_{\mathcal{X}} V$는
열린 몰입이다. 이 도식과 보조정리의 명제에 등장하는 여러 차원의 정의를
함께 사용하여, 보조정리의 확인을 알려진 스킴의 경우로 환원할 것이다.

\medskip\noindent
$w \in |W|$를 고정하고, 그 상을 $w' \in |\mathcal{W}|$,
그 상을 $t \in |\mathcal{T}|$, $|V|$에서의 상을 $v$,
$|\mathcal{X}|$에서의 상을 $x$라 하자.
사실상 정의에 의해($\mathcal{W}$가
$\mathcal{T}\times_{\mathcal{X}} V$에서 열려 있고 밑변환의 섬유가
섬유의 밑변환이라는 사실을 사용하여) 등식
$$
\dim_v V_x = \dim_{w'} \mathcal{W}_t
$$
및
$$
\dim_t \mathcal{T}_x = \dim_{w'} \mathcal{W}_v.
$$
을 얻는다. 보조정리
\ref{stacks-geometry-lemma-behaviour-of-dimensions-wrt-smooth-morphisms}에 의해
(우리 도식의 대각 화살표와 오른쪽 수직 화살표는 $W$와 $V$를 각각
스택 $\mathcal{T}$와 $\mathcal{X}$의 스킴에 의한 매끄러운 덮개로 실현한다)
$$
\dim_t \mathcal{T} = \dim_w W - \dim_w W_t
$$
이고
$$
\dim_x \mathcal{X} = \dim_v V - \dim_v V_x.
$$
이다. 이 등식들을 합치면
$$
\dim_t \mathcal{T}_x - \dim_t \mathcal{T} + \dim_x \mathcal{X}
= \dim_{w'} \mathcal{W}_v - \dim_w W + \dim_w W_t + \dim_v V -
\dim_{w'} \mathcal{W}_t
$$
를 얻는다. $W \to \mathcal{W}$는 매끄러운 전사이므로,
($W \to \mathcal{W}$를 $V$ 위의 사상으로 생각하여)
사상 $\Spec \kappa(v) \to V$ 위로
밑변환해도 마찬가지이다. 이 매끄러운 사상에서 다음 두 등식 중
첫 번째를 얻는다.
$$
\dim_w W_v - \dim_{w'} \mathcal{W}_v = \dim_w (W_v)_{w'} = \dim_w W_{w'};
$$
두 번째 등식은 등장하는 두 섬유를 직접 비교하면 따른다.
마찬가지로 $W \to \mathcal{W}$를 $\mathcal{T}$ 위의 스킴들의 사상으로
생각하고 점 $t \in |\mathcal{T}|$의 어떤 대표 위로 밑변환하면 등식
$$
\dim_w W_t - \dim_{w'} \mathcal{W}_t = \dim_w (W_t)_{w'} = \dim_w W_{w'}.
$$
을 얻는다. 모든 것을 합치면
$$
\dim_t \mathcal{T}_x - \dim_t \mathcal{T} + \dim_x \mathcal{X}
=  \dim_w W_v - \dim_w W + \dim_v V.
$$
을 얻는다. 우리의 목표는 이 등식의 좌변이 공집합이 아닌
$t$의 열린부분집합에서 사라짐을 보이는 것이다.
$w$가 $W$의 공집합이 아닌 열린부분집합 위에서 변하면 그 상
$t \in |\mathcal{T}|$도 $|\mathcal{T}|$의 공집합이 아닌
열린부분집합 위에서 변한다($W \to \mathcal{T}$가 매끄럽기 때문이다).

\medskip\noindent
따라서 다음을 보이는 것으로 환원된다. $W\to V$가 기약인 국소 뇌터
스킴들의 국소적으로 유한형이고 스킴론적으로 우세한 사상이면,
$\dim_w W_v =\dim_w W - \dim_v V$가 성립하는 점 $w\in W$들의
공집합이 아닌 열린부분집합이 존재한다(여기서 $v$는 $V$에서 $w$의 상이다).
이는 표준적인 사실이며, 독자의 편의를 위해 그 증명을 상기한다.

\medskip\noindent
이 등식의 성립 여부를 바꾸지 않고 $W$와 $V$를 그 바탕 축소
부분스킴으로 대체할 수 있다. 따라서 이들이 실제로 정수적 스킴이라고
가정해도 좋다. $\dim_w W_v$는 $W$ 위에서 국소 상수이므로,
필요하면 $W$를 공집합이 아닌 열린부분집합으로 대체하여
$\dim_w W_v$가 어떤 $d$와 같은 상수라고 가정해도 좋다.
이 열린부분집합을 아핀으로 택하면 사상 $W\to V$가 실제로 유한형이라고도
가정할 수 있다. 필요하면 $V$를 공집합이 아닌 열린부분집합으로 대체하고
(그런 다음 이 열린부분집합 위로 $W$를 당긴다. 준콤팩트이고
스킴론적으로 우세한 사상의 평탄 밑변환은 여전히 스킴론적으로
우세하므로, 그 당김은 공집합이 아니다), 더 나아가 $W$가 $V$ 위에서
평탄하다고 가정할 수 있다. 따라서 『스킴의 사상』의 정의
\ref{morphisms-definition-relative-dimension-d}의 의미에서
사상 $W\to V$는 상대 차원 $d$이고, 『스킴의 사상』의 보조정리
\ref{morphisms-lemma-rel-dimension-dimension}에 의해 요구한 대로
$\dim_w(W) = \dim_v(V) + d$이다.
\end{proof}

\begin{remark}
\label{stacks-geometry-remark-negative-dimension}
앞의 보조정리의 맥락에서는 $\dim \mathcal{T} \geq \dim \mathcal{Z}$일
필요가 없음을 언급한다. 이는 보조정리 명제의 부등식과 모순되지 않는다.
사상 $f$의 섬유도 대수적 스택이므로 음의 차원을 가질 수 있기 때문이다.
$k$를 체라 하고 보조정리를 사상
$[\Spec k/\mathbf{G}_m] \to \Spec k$에 적용하면 이를 볼 수 있다.

\medskip\noindent
보조정리 명제의 사상 $f$가 준-DM이라고 가정하면
(『스택의 사상』의 정의
\ref{stacks-morphisms-definition-separated}의 의미에서이다. 예를 들어
대수공간으로 표현가능한 사상은 준-DM이다), 표적의 점 위에서 사상의
섬유는 준-DM 대수적 스택이므로 차원이 음이 아니다. 이 경우 보조정리는
실제로 $\dim \mathcal{T} \geq \dim \mathcal{Z}$임을 함의한다.
사실 다음과 같은 더 일반적인 결과를 얻는다.
\end{remark}

\begin{lemma}
\label{stacks-geometry-lemma-dims-of-images-two}
$f: \mathcal{T} \to \mathcal{X}$를 제이컵슨이고 유사 코테너리이며
국소 뇌터인 대수적 스택들의 국소적으로 유한형인 준-DM 사상이라 하자.
원천은 기약이고 표적은 준분리라고 하며,
$\mathcal{Z} \hookrightarrow \mathcal{X}$를 $\mathcal{T}$의
스킴론적 상이라 하자. 그러면
$\dim \mathcal{Z} \leq \dim \mathcal{T}$이고, 더 나아가 다음 두 조건 중
정확히 하나가 성립한다.
\begin{enumerate}
\item 모든 유한형 점 $t \in |T|$에 대해
$\dim_t(\mathcal{T}_{f(t)}) > 0$이고, 이 경우
$\dim \mathcal{Z} < \dim \mathcal{T}$이다. 또는
\item $\mathcal{T}$와 $\mathcal{Z}$의 차원이 같다.
\end{enumerate}
\end{lemma}

\begin{proof}
앞의 주석에서 관찰했듯이 준-DM 스택의 차원은 항상 음이 아니다.
따라서 모든 $t \in |\mathcal{T}|$에 대해
$\dim_t \mathcal{T}_{f(t)} \geq 0$이고, 점들
$t\in |\mathcal{T}|$의 조밀한 열린부분집합에서 등식
$$
\dim_t \mathcal{T}_{f(t)} = \dim_t \mathcal{T} - \dim_{f(t)} \mathcal{Z}
$$
이 성립한다고 결론 내린다.
\end{proof}





\section{국소환의 차원}
\label{stacks-geometry-section-dimension-local-ring}

\noindent
대수적 스택은 통상적인 의미의 국소환을 실제로 갖지는 않지만,
국소환의 차원을 다음과 같이 정의할 수 있다.

\begin{lemma}
\label{stacks-geometry-lemma-dimension-local-ring-pre}
$\mathcal{X}$를 국소 뇌터 대수적 스택이라 하자.
$U \to \mathcal{X}$를 매끄러운 사상이라 하고 $u \in U$라 하자. 그러면
$$
\dim(\mathcal{O}_{U, \overline{u}}) -
\dim(\mathcal{O}_{R_u, e(\overline{u})}) =
2\dim(\mathcal{O}_{U, \overline{u}}) -
\dim(\mathcal{O}_{R, e(\overline{u})})
$$
이다. 여기서 $R = U \times_\mathcal{X} U$이고 사영은
$s, t : R \to U$, 대각사상은 $e : U \to R$이며,
$R_u$는 $u$ 위에서 $s : R \to U$의 섬유이다.
\end{lemma}

\begin{proof}
이는 $s : \mathcal{O}_{U, \overline{u}} \to
\mathcal{O}_{R, e(\overline{u})}$가 뇌터 국소환들의 평탄 국소
준동형사상이어서 『대수학』의 보조정리
\ref{algebra-lemma-dimension-base-fibre-equals-total}에 의해
$$
\dim(\mathcal{O}_{R, e(\overline{u})}) =
\dim(\mathcal{O}_{U, \overline{u}}) +
\dim(\mathcal{O}_{R_u, e(\overline{u})})
$$
이기 때문이다.
\end{proof}

\begin{lemma}
\label{stacks-geometry-lemma-dimension-local-ring}
$\mathcal{X}$를 국소 뇌터 대수적 스택이라 하자.
$x \in |\mathcal{X}|$를 유한형 점이라 하자
『스택의 사상』의 정의
\ref{stacks-morphisms-definition-finite-type-point}).
$d \in \mathbf{Z}$라 하자. 다음은 동치이다.
\begin{enumerate}
\item 스킴 $U$, 매끄러운 사상 $U \to \mathcal{X}$,
$x$로 사상되는 유한형 점 $u \in U$가 존재하여
$2\dim(\mathcal{O}_{U, \overline{u}}) -
\dim(\mathcal{O}_{R, e(\overline{u})}) = d$이다. 그리고
\item 임의의 스킴 $U$, 매끄러운 사상 $U \to \mathcal{X}$,
$x$로 사상되는 유한형 점 $u \in U$에 대해
$2\dim(\mathcal{O}_{U, \overline{u}}) -
\dim(\mathcal{O}_{R, e(\overline{u})}) = d$이다.
\end{enumerate}
여기서 $R = U \times_\mathcal{X} U$이고 사영은
$s, t : R \to U$, 대각사상은 $e : U \to R$이며,
$R_u$는 $u$ 위에서 $s : R \to U$의 섬유이다.
\end{lemma}

\begin{proof}
$(U, u)$와 $(U', u')$를 $x$의 두 매끄러운 근방이라 하고,
$u$와 $u'$을 유한형 점이라 하자. $U$와 $U'$을 축소한 뒤
$u$와 $u'$이 닫힌점이라고 가정해도 좋다(유한형 점의 정의에 의해).
그런 다음 전사 에탈 사상 $W \to U \times_\mathcal{X} U'$을 택한다.
$W_u$를 $u$ 위에서 $W \to U$의 섬유라 하고,
$W_{u'}$을 $u'$ 위에서 $W \to U'$의 섬유라 하자.
$u$와 $u'$이 $|\mathcal{X}|$의 같은 점으로 사상되므로
$W_u \cap W_{u'}$는 공집합이 아니다.
따라서 $u$와 $u'$ 모두로 사상되는 닫힌점 $w \in W$를 택할 수 있다.
이로써 다음 문단의 논의로 환원된다.

\medskip\noindent
$(U', u') \to (U, u)$가 $x$의 매끄러운 근방들의 매끄러운 사상이고,
$u$와 $u'$은 닫힌점이라고 가정하자. 목표는 $(U, u)$에 대해 정의한
불변량과 $(U', u')$에 대해 정의한 불변량이 같음을 증명하는 것이다.
이를 보기 위해 $\mathcal{O}_{U, u} \to \mathcal{O}_{U', u'}$가
뇌터 국소환들의 평탄 국소 준동형사상임을 관찰하자. 따라서
『대수학』의 보조정리
\ref{algebra-lemma-dimension-base-fibre-equals-total}에 의해
$$
\dim(\mathcal{O}_{U', \overline{u}'}) =
\dim(\mathcal{O}_{U, \overline{u}}) +
\dim(\mathcal{O}_{U'_u, \overline{u}'})
$$
이다. (국소환과 그 엄밀 헨젤화의 성질을 연결하는 모든 단계를 다루는
일은 생략한다. 『대수학 더 보기』의 절
\ref{more-algebra-section-permanence-henselization}을 보라.)
한편
$$
R' = U' \times_{U, t} R \times_{s, U} U'
$$
이므로
$$
\dim(\mathcal{O}_{R', e(\overline{u}')}) =
\dim(\mathcal{O}_{R, e(\overline{u})}) +
\dim(\mathcal{O}_{U'_u \times_u U'_u, (\overline{u}', \overline{u}')})
$$
임을 알 수 있다. 보조정리를 증명하려면
$$
\dim(\mathcal{O}_{U'_u \times_u U'_u, (\overline{u}', \overline{u}')}) =
2\dim(\mathcal{O}_{U'_u, \overline{u}'})
$$
임을 보이면 충분하다. 이것은 항상 참은 아니다
(예를 들어 $U'_u$가 곡선이고 $u'$이 이 곡선의 일반점인 경우).
그러나 $u'$은 $u$ 위에서 국소적으로 유한형인 대수공간 $U'_u$의
닫힌점임을 안다. 이 경우 등식은 다음 두 사실 때문에 성립한다.
첫째, 『다양체』의 보조정리
\ref{varieties-lemma-dimension-product-locally-algebraic}에 의해
$\dim_{(u', u')}(U'_u \times_u U'_u) = 2\dim_{u'}(U'_u)$이다.
둘째, 국소적으로 대수적인 스킴의 닫힌점에서는 차원과 국소환의 차원이
일치한다. 『다양체』의 보조정리
\ref{varieties-lemma-dimension-locally-algebraic}를 보라.
스킴에 관한 이 결과들을 대수공간의 언어로 옮기는 일은 생략한다.
\end{proof}

\begin{definition}
\label{stacks-geometry-definition-dimension-local-ring}
$\mathcal{X}$를 국소 뇌터 대수적 스택이라 하자.
$x \in |\mathcal{X}|$를 유한형 점이라 하자.
보조정리 \ref{stacks-geometry-lemma-dimension-local-ring}의 동치인 조건들이 만족될 때,
{\it $x$에서 $\mathcal{X}$의 국소환의 차원}은 $d \in \mathbf{Z}$이다.
\end{definition}

\noindent
분명 이 정의의 동기는 보조정리 \ref{stacks-geometry-lemma-dimension-local-ring-pre}와
『스택의 성질』의 정의
\ref{stacks-properties-definition-dimension-at-point}이다.
이 절을 마치며, $\mathcal{X}$의 $x$에서의 버설 환의 성질로
$\dim_x(\mathcal{X})$를 계산할 수 있게 하는 공식을 확립한다.

\begin{lemma}
\label{stacks-geometry-lemma-dimension-formula}
$\mathcal{X}$가 국소 뇌터 스킴 $S$ 위에서 국소적으로 유한형인
대수적 스택이라고 하자. $x_0 : \Spec(k) \to \mathcal{X}$를 사상이라
하자. 여기서 $k$는 $S$ 위의 유한형 체이다.
주석 \ref{stacks-geometry-remark-groupoid-defo}에서와 같이
$\mathcal{F}_{\mathcal{X}, k, x_0}$를 잉여체가 $k$인 뇌터 완비 국소
$S$-대수들의 쌍대준군 $(A, B, s, t, c)$로 나타내자. 그러면
$$
\text{국소환 차원: }\mathcal{X}\text{의 점 }x_0 =
2\dim A - \dim B
$$
이다.
\end{lemma}

\begin{proof}
$s \in S$를 $x_0$의 상이라 하자. $\mathcal{O}_{S, s}$가 G-환이면
(실제로 거의 항상 만족되는 조건이다) 다음과 같이 보조정리를 증명할 수 있다.
보조정리 \ref{stacks-geometry-lemma-Artin-approximation-by-smooth-morphism}에 의해,
원천이 스킴이고 잉여체가 $k$인 점 $u_0 \in U$를 포함하는
매끄러운 사상 $U \to \mathcal{X}$를 찾을 수 있다. 이때 유도된 사상
$\Spec(k) \to U \to \mathcal{X}$는 $x_0$와 일치하고
$A = \mathcal{O}_{U, u_0}^\wedge$이다.
$R = U \times_\mathcal{X} U$로 쓰자. 그러면
$\mathcal{O}_{R, e(u_0)}^\wedge$를 $B$와 동일시할 수 있다.
따라서 등식은 정의에서 따른다.

\medskip\noindent
이 증명의 나머지에서는 보조정리를 일반적으로 증명하는 방법을 설명하지만,
독자에게는 이 부분을 건너뛰기를 권한다.

\medskip\noindent
먼저 우변이 $(A, B, s, t, c)$의 선택에 무관함을 보이자.
$(A', B', s', t', c')$를 두 번째 선택이라고 하자.
$A$와 $A'$은 $\mathcal{X}$의 $x_0$에서의 버설 환이므로,
필요하면 $A$와 $A'$을 맞바꾼 뒤 $A$와 $A'$ 위의 주어진 버설 형식적 대상
$\xi$와 $\xi'$에 양립하는 형식적으로 매끄러운 사상
$A \to A'$을 택할 수 있다.
$\widehat{\mathcal{C}}_\Lambda$가 쌍대곱을 가지며 이들이 $\Lambda$ 위의
완비 텐서곱으로 주어진다는 사실을 상기하자. 『형식적 변형 이론』의
보조정리 \ref{formal-defos-lemma-CLambdahat-coproducts}를 보라.
그러면 $B$는 $\xi$를 $A \widehat{\otimes}_\Lambda A$로 보내어 얻은
두 대상 사이의 동형사상 함자를 프로표현한다. $B'$도 마찬가지이다.
따라서
$$
B' =
B \otimes_{(A \widehat{\otimes}_\Lambda A)}
(A' \widehat{\otimes}_\Lambda A')
$$
라고 결론 내린다. 다음 사상이 형식적으로 매끄럽고, 그 상대 차원이
형식적으로 매끄러운 사상 $A \to A'$의 상대 차원의 $2$배와 같음을
곧바로 알 수 있다.
$$
A \widehat{\otimes}_\Lambda A \longrightarrow
A \widehat{\otimes}_\Lambda A' \longrightarrow
A' \widehat{\otimes}_\Lambda A'
$$
(이는 일반 원리에서 따르지만, 이 특별한 경우에는 $A'$이 $r$개 변수에
대한 $A$ 위의 멱급수환이기 때문이기도 하다.) 따라서 원하는 대로
$B \to B'$은 상대 차원이 $2(\dim(A') - \dim(A))$인
형식적으로 매끄러운 사상이다.

\medskip\noindent
다음으로 $l/k$를 유한 확대라 하자.
$x_{l, 0} : \Spec(l) \to \mathcal{X}$를 유도된 점이라 하자.
공식의 우변이 $x_0$에 대해서나 $x_{l, 0}$에 대해서나 같다고 주장한다.
이는 보조정리 \ref{stacks-geometry-lemma-versal-ring-field-extension}에서와 같은
$A \to A'$을 택하고 앞 문단과 정확히 똑같이 논증하면 보일 수 있다.
세부사항은 생략한다.

\medskip\noindent
마지막으로 보조정리 \ref{stacks-geometry-lemma-versal-ring-flat}의 증명에서와 같이
논증하면, 앞의 두 문단의 양립성을 사용하여 $A$가 $\mathcal{X}$ 위에서
매끄러운 어떤 스킴 $U$의 $u_0$에서의 완비 국소환이고 $u_0$가 유한형 점인 경우
(첫 문단에서 논의한 경우)로 환원할 수 있다. 세부사항은 생략한다.
\end{proof}










\noindent
이제 다음 정의가 의미 있음을 보이기에 충분한 준비가 되었다.
